% Options for packages loaded elsewhere
\PassOptionsToPackage{unicode}{hyperref}
\PassOptionsToPackage{hyphens}{url}
\documentclass[
  french,
]{article}
\usepackage{xcolor}
\usepackage[margin=2cm]{geometry}
\usepackage{amsmath,amssymb}
\setcounter{secnumdepth}{-\maxdimen} % remove section numbering
\usepackage{iftex}
\ifPDFTeX
  \usepackage[T1]{fontenc}
  \usepackage[utf8]{inputenc}
  \usepackage{textcomp} % provide euro and other symbols
\else % if luatex or xetex
  \usepackage{unicode-math} % this also loads fontspec
  \defaultfontfeatures{Scale=MatchLowercase}
  \defaultfontfeatures[\rmfamily]{Ligatures=TeX,Scale=1}
\fi
\usepackage{lmodern}
\ifPDFTeX\else
  % xetex/luatex font selection
\fi
% Use upquote if available, for straight quotes in verbatim environments
\IfFileExists{upquote.sty}{\usepackage{upquote}}{}
\IfFileExists{microtype.sty}{% use microtype if available
  \usepackage[]{microtype}
  \UseMicrotypeSet[protrusion]{basicmath} % disable protrusion for tt fonts
}{}
\makeatletter
\@ifundefined{KOMAClassName}{% if non-KOMA class
  \IfFileExists{parskip.sty}{%
    \usepackage{parskip}
  }{% else
    \setlength{\parindent}{0pt}
    \setlength{\parskip}{6pt plus 2pt minus 1pt}}
}{% if KOMA class
  \KOMAoptions{parskip=half}}
\makeatother
\usepackage{color}
\usepackage{fancyvrb}
\newcommand{\VerbBar}{|}
\newcommand{\VERB}{\Verb[commandchars=\\\{\}]}
\DefineVerbatimEnvironment{Highlighting}{Verbatim}{commandchars=\\\{\}}
% Add ',fontsize=\small' for more characters per line
\newenvironment{Shaded}{}{}
\newcommand{\AlertTok}[1]{\textcolor[rgb]{1.00,0.00,0.00}{\textbf{#1}}}
\newcommand{\AnnotationTok}[1]{\textcolor[rgb]{0.38,0.63,0.69}{\textbf{\textit{#1}}}}
\newcommand{\AttributeTok}[1]{\textcolor[rgb]{0.49,0.56,0.16}{#1}}
\newcommand{\BaseNTok}[1]{\textcolor[rgb]{0.25,0.63,0.44}{#1}}
\newcommand{\BuiltInTok}[1]{\textcolor[rgb]{0.00,0.50,0.00}{#1}}
\newcommand{\CharTok}[1]{\textcolor[rgb]{0.25,0.44,0.63}{#1}}
\newcommand{\CommentTok}[1]{\textcolor[rgb]{0.38,0.63,0.69}{\textit{#1}}}
\newcommand{\CommentVarTok}[1]{\textcolor[rgb]{0.38,0.63,0.69}{\textbf{\textit{#1}}}}
\newcommand{\ConstantTok}[1]{\textcolor[rgb]{0.53,0.00,0.00}{#1}}
\newcommand{\ControlFlowTok}[1]{\textcolor[rgb]{0.00,0.44,0.13}{\textbf{#1}}}
\newcommand{\DataTypeTok}[1]{\textcolor[rgb]{0.56,0.13,0.00}{#1}}
\newcommand{\DecValTok}[1]{\textcolor[rgb]{0.25,0.63,0.44}{#1}}
\newcommand{\DocumentationTok}[1]{\textcolor[rgb]{0.73,0.13,0.13}{\textit{#1}}}
\newcommand{\ErrorTok}[1]{\textcolor[rgb]{1.00,0.00,0.00}{\textbf{#1}}}
\newcommand{\ExtensionTok}[1]{#1}
\newcommand{\FloatTok}[1]{\textcolor[rgb]{0.25,0.63,0.44}{#1}}
\newcommand{\FunctionTok}[1]{\textcolor[rgb]{0.02,0.16,0.49}{#1}}
\newcommand{\ImportTok}[1]{\textcolor[rgb]{0.00,0.50,0.00}{\textbf{#1}}}
\newcommand{\InformationTok}[1]{\textcolor[rgb]{0.38,0.63,0.69}{\textbf{\textit{#1}}}}
\newcommand{\KeywordTok}[1]{\textcolor[rgb]{0.00,0.44,0.13}{\textbf{#1}}}
\newcommand{\NormalTok}[1]{#1}
\newcommand{\OperatorTok}[1]{\textcolor[rgb]{0.40,0.40,0.40}{#1}}
\newcommand{\OtherTok}[1]{\textcolor[rgb]{0.00,0.44,0.13}{#1}}
\newcommand{\PreprocessorTok}[1]{\textcolor[rgb]{0.74,0.48,0.00}{#1}}
\newcommand{\RegionMarkerTok}[1]{#1}
\newcommand{\SpecialCharTok}[1]{\textcolor[rgb]{0.25,0.44,0.63}{#1}}
\newcommand{\SpecialStringTok}[1]{\textcolor[rgb]{0.73,0.40,0.53}{#1}}
\newcommand{\StringTok}[1]{\textcolor[rgb]{0.25,0.44,0.63}{#1}}
\newcommand{\VariableTok}[1]{\textcolor[rgb]{0.10,0.09,0.49}{#1}}
\newcommand{\VerbatimStringTok}[1]{\textcolor[rgb]{0.25,0.44,0.63}{#1}}
\newcommand{\WarningTok}[1]{\textcolor[rgb]{0.38,0.63,0.69}{\textbf{\textit{#1}}}}
\usepackage{longtable,booktabs,array}
\usepackage{caption}
\captionsetup[table]{skip=6pt}
\newcounter{none} % for unnumbered tables
\usepackage{calc} % for calculating minipage widths
% Correct order of tables after \paragraph or \subparagraph
\usepackage{etoolbox}
\makeatletter
\patchcmd\longtable{\par}{\if@noskipsec\mbox{}\fi\par}{}{}
\makeatother
\ifLuaTeX
\usepackage[bidi=basic,shorthands=off]{babel}
\else
\usepackage[bidi=default,shorthands=off]{babel}
\fi
\ifLuaTeX
  \usepackage{selnolig} % disable illegal ligatures
\fi
\setlength{\emergencystretch}{3em} % prevent overfull lines
\providecommand{\tightlist}{%
  \setlength{\itemsep}{0pt}\setlength{\parskip}{0pt}}
\usepackage{bookmark}
\IfFileExists{xurl.sty}{\usepackage{xurl}}{} % add URL line breaks if available
\urlstyle{same}
% fallback for those not using the hyperref driver hyperxmp:
\makeatletter
\@ifundefined{xmpquote}{\newcommand{\xmpquote}[1]{#1}}{}
\makeatother
\hypersetup{
  pdflang={fr},
  hidelinks,
  pdfcreator={LaTeX via pandoc}}

\author{}
\date{}

\begin{document}

\section{Manuel de Référence Utilisateur et Guide Pratique - Solveur
COVOF
2D}\label{manuel-de-ruxe9fuxe9rence-utilisateur-et-guide-pratique---solveur-covof-2d}

\subsection{Guide Complet d'Installation, Paramétrage, Exécution et
Post-Traitement}\label{guide-complet-dinstallation-paramuxe9trage-exuxe9cution-et-post-traitement}

\subsubsection{Version 2.1.0 --- Conforme aux Standards du Laboratoire
TREFLE / I2M (Référence
Académique)}\label{version-2.1.0-conforme-aux-standards-du-laboratoire-trefle-i2m-ruxe9fuxe9rence-acaduxe9mique}

\begin{center}\rule{0.5\linewidth}{0.5pt}\end{center}

\subsection{Sommaire Général}\label{sommaire-guxe9nuxe9ral}

\begin{enumerate}
\def\labelenumi{\arabic{enumi}.}
\tightlist
\item
  \hyperref[1-pruxe9sentation-et-fonctionnalituxe9s-du-logiciel]{Présentation
  et Fonctionnalités du Logiciel}

  \begin{itemize}
  \tightlist
  \item
    1.1. Philosophie du Solveur COVOF 2D
  \item
    1.2. Matrice Globale des Fonctionnalités Physiques et Numériques
  \end{itemize}
\item
  \hyperref[2-installation-duxe9pendances-et-compilation]{Installation,
  Dépendances et Compilation}

  \begin{itemize}
  \tightlist
  \item
    2.1. Prérequis Système et Compilateurs C++20
  \item
    2.2. Bibliothèques Externes (MPI, Python, Outils de Tracé)
  \item
    2.3. Procédure de Compilation pas-à-pas avec CMake
  \item
    2.4. Arborescence et Organisation des Fichiers du Projet
  \end{itemize}
\item
  \hyperref[3-format-complet-des-fichiers-de-configuration-json]{Format
  Complet des Fichiers de Configuration (JSON)}

  \begin{itemize}
  \tightlist
  \item
    3.1. Structure Générale d'un Cas Test
  \item
    3.2. Dictionnaire Exhaustif des Modules et Mots-Clés

    \begin{itemize}
    \tightlist
    \item
      3.2.1. Module \texttt{mesh} (Maillage et Dimensions Géométriques)
    \item
      3.2.2. Module \texttt{physics} (Propriétés Physiques et Modèles
      d'Écoulement)
    \item
      3.2.3. Module \texttt{thermal} (Couplage Énergétique et
      Boussinesq)
    \item
      3.2.4. Module \texttt{time} (Intégration Temporelle et Contrôle
      CFL)
    \item
      3.2.5. Module \texttt{boundary\_conditions} (Conditions aux
      Limites Fluides et Thermiques)
    \item
      3.2.6. Module \texttt{immersed\_boundary} (Frontières Immergées et
      Pénalisation Brinkman)
    \item
      3.2.7. Module \texttt{initial\_condition} (Champs Initiaux et
      Géométries d'Interface)
    \item
      3.2.8. Module \texttt{solver} (Méthode Numérique, Solveur Krylov
      et Préconditionneur)
    \item
      3.2.9. Module \texttt{output} (Sorties Graphiques VTK et
      Diagnostics CSV)
    \end{itemize}
  \end{itemize}
\item
  \hyperref[4-tutoriel-pas-uxe0-pas--construire-un-cas-de-simulation]{Tutoriel
  pas-à-pas : Construire un Cas de Simulation}

  \begin{itemize}
  \tightlist
  \item
    4.1. Tutoriel 1 : Écoulement Monophasique dans une Cavité Entraînée
    (\(Re=1000\))
  \item
    4.2. Tutoriel 2 : Convection Naturelle Thermique de Rayleigh-Bénard
    / Vahl Davis
  \item
    4.3. Tutoriel 3 : Écoulement Diphasique à Surface Libre (Rupture de
    Barrage)
  \end{itemize}
\item
  \hyperref[5-exuxe9cution-du-code-en-suxe9quentiel-et-paralluxe8le-openmpi]{Exécution
  du Code en Séquentiel et Parallèle (OpenMPI)}

  \begin{itemize}
  \tightlist
  \item
    5.1. Lancement Local sur Poste de Travail ou Ordinateur Portable
  \item
    5.2. Lancement sur Supercalculateur et Gestionnaire de Tâches SLURM
  \item
    5.3. Décomposition de Domaine MPI et Équilibrage de Charge
  \end{itemize}
\item
  \hyperref[6-suivi-dexuxe9cution-et-diagnostic-des-calculs-en-temps-ruxe9el]{Suivi
  d'Exécution et Diagnostic des Calculs en Temps Réel}

  \begin{itemize}
  \tightlist
  \item
    6.1. Interprétation du Tableau de Bord Console
  \item
    6.2. Détection Précoce des Instabilités et Divergences
  \item
    6.3. Guide de Dépannage et Stratégies de Remédiation
  \end{itemize}
\item
  \hyperref[7-post-traitement-et-analyse-graphique-avec-paraview]{Post-Traitement
  et Analyse Graphique avec ParaView}

  \begin{itemize}
  \tightlist
  \item
    7.1. Chargement des Fichiers \texttt{.pvtr} et \texttt{.vtr}
  \item
    7.2. Extraction des Lignes de Courant et Profils de Vitesse
  \item
    7.3. Tracé des Isothermes et Isosurfaces d'Interface (\(C = 0.5\))
  \item
    7.4. Calcul de la Vorticité et des Gradients de Pression
  \end{itemize}
\item
  \hyperref[8-automatisation-des-benchmarks-et-scripts-de-vuxe9rification]{Automatisation
  des Benchmarks et Scripts de Vérification}

  \begin{itemize}
  \tightlist
  \item
    8.1. Le Script d'Orchestration Globale
    \texttt{run\_all\_benchmarks.sh}
  \item
    8.2. Vérification Statistique Automatisée \texttt{verify\_cases.py}
  \item
    8.3. Production Automatique des Figures de Rapport
    \texttt{generate\_validation\_figures.py}
  \end{itemize}
\item
  \hyperref[9-foire-aux-questions-faq-et-astuces-pratiques]{Foire Aux
  Questions (FAQ) et Astuces Pratiques}
\end{enumerate}

\begin{center}\rule{0.5\linewidth}{0.5pt}\end{center}

\subsection{1. Présentation et Fonctionnalités du
Logiciel}\label{pruxe9sentation-et-fonctionnalituxe9s-du-logiciel}

\subsubsection{1.1. Philosophie du Solveur COVOF
2D}\label{philosophie-du-solveur-covof-2d}

Le code \textbf{COVOF 2D} est un environnement complet de simulation
numérique des écoulements de fluides incompressibles en deux dimensions
d'espace. Il a été conçu avec les impératifs d'ingénierie logicielle
suivants : - \textbf{Fidélité physique maximale} : Implémentation
stricte des équations de Navier-Stokes à viscosité et masse volumique
variables, de l'approximation de Boussinesq pour la thermique, du modèle
CSF de Brackbill pour les interfaces et de la pénalisation volumique de
Brinkman pour les obstacles. - \textbf{Robustesse algorithmique absolue}
: Utilisation de la méthode du \textbf{Lagrangien Augmenté},
garantissant l'incompressibilité sans instabilités de pression et
éliminant le besoin de résoudre une équation de Poisson mal
conditionnée. - \textbf{Hautes performances et montée en charge} :
Parallélisme natif via \textbf{OpenMPI}, décomposition de domaine
cartésienne 2D équilibrée, communications non bloquantes, structures de
données à accès mémoire contigu sans aucune allocation dynamique pendant
la boucle en temps. - \textbf{Simplicité d'utilisation} : Fichiers
d'entrée au format standard \textbf{JSON}, entièrement autodocumentés,
sans compilation préalable pour changer de géométrie, de maillage ou de
conditions physiques. - \textbf{Interfaçage standardisé} : Sorties
graphiques directes au format XML VTK multi-blocs (\texttt{.pvtr} et
\texttt{.vtr}), visualisables immédiatement sous ParaView ou VisIt sans
utilitaire de conversion.

\subsubsection{1.2. Matrice Globale des Fonctionnalités Physiques et
Numériques}\label{matrice-globale-des-fonctionnalituxe9s-physiques-et-numuxe9riques}

{\def\LTcaptype{none} % do not increment counter
\begin{longtable}[]{@{}
  >{\raggedright\arraybackslash}p{(\linewidth - 8\tabcolsep) * \real{0.1818}}
  >{\raggedright\arraybackslash}p{(\linewidth - 8\tabcolsep) * \real{0.1818}}
  >{\centering\arraybackslash}p{(\linewidth - 8\tabcolsep) * \real{0.2273}}
  >{\centering\arraybackslash}p{(\linewidth - 8\tabcolsep) * \real{0.2273}}
  >{\raggedright\arraybackslash}p{(\linewidth - 8\tabcolsep) * \real{0.1818}}@{}}
\toprule\noalign{}
\begin{minipage}[b]{\linewidth}\raggedright
Domaine
\end{minipage} & \begin{minipage}[b]{\linewidth}\raggedright
Fonctionnalité Disponible
\end{minipage} & \begin{minipage}[b]{\linewidth}\centering
Séquentiel
\end{minipage} & \begin{minipage}[b]{\linewidth}\centering
Parallèle (OpenMPI)
\end{minipage} & \begin{minipage}[b]{\linewidth}\raggedright
Détails Techniques \& Référence
\end{minipage} \\
\midrule\noalign{}
\endhead
\bottomrule\noalign{}
\endlastfoot
\textbf{Écoulement} & Monophasique incompressible & ✅ & ✅ &
Navier-Stokes avec tenseur visqueux complet \\
& Viscosité variable continue & ✅ & ✅ & Modèle exponentiel ou
analytique nodale \\
& Thermoconvection naturelle & ✅ & ✅ & Boussinesq, nombre de Nusselt
automatisé \\
& Diphasique non miscible & ✅ & ✅ & VOF (Volume of Fluid)
conservatif \\
& Tension superficielle & ✅ & ✅ & Modèle CSF Brackbill
Balanced-Force \\
& Frontières immergées & ✅ & ✅ & Pénalisation Brinkman (marche,
cylindre, L) \\
\textbf{Maillage} & Cartésien MAC 2D décalé & ✅ & ✅ & Staggered grid
de Harlow \& Welch (1965) \\
& Périodicité & ✅ & ✅ & Raccord torique d'indices et échanges de
halos \\
\textbf{Solveurs} & BiCGSTAB(2) & ✅ & ✅ & Solveur Krylov résiduel
d'ordre 2 (Sleijpen) \\
& BiCGSTAB standard & ✅ & ✅ & Gradient biconjugué stabilisé \\
& Gradient Conjugué (PCG) & ✅ & ✅ & Dédié aux matrices symétriques
définies positives \\
& GMRES(m) & ✅ & ✅ & Redémarrages avec rotations de Givens \\
& MINRES & ✅ & ✅ & Pour systèmes symétriques indéfinis \\
& MUMPS / PETSc / HYPRE & Option & Option & Interfaces prêtes pour
solveurs externes directs/AMG \\
\textbf{Préconditionnement} & Jacobi (Diagonal) & ✅ & ✅ & Parallélisme
parfait sans communication inter-cœurs \\
& SSOR / ILU(0) & ✅ & ✅ & Factorisation incomplète par sous-domaine \\
\end{longtable}
}

\begin{center}\rule{0.5\linewidth}{0.5pt}\end{center}

\subsection{2. Installation, Dépendances et
Compilation}\label{installation-duxe9pendances-et-compilation}

\subsubsection{2.1. Prérequis Système et Compilateurs
C++20}\label{pruxe9requis-systuxe8me-et-compilateurs-c20}

COVOF 2D exploite les fonctionnalités avancées du standard
\textbf{C++20} (concepts, vues, attributs d'optimisation). -
\textbf{Systèmes d'exploitation supportés} : Linux (Ubuntu 20.04+, RHEL
8+, Debian 11+), macOS (macOS 12+ avec Clang/AppleClang ou GCC
Homebrew). - \textbf{Compilateurs C++ certifiés} : - \texttt{GCC}
version 10.3 ou ultérieure (version 11+ recommandée). -
\texttt{Clang\ /\ LLVM} version 13.0 ou ultérieure. -
\texttt{AppleClang} version 13.0 ou ultérieure sous macOS.

\subsubsection{2.2. Bibliothèques
Externes}\label{bibliothuxe8ques-externes}

\begin{enumerate}
\def\labelenumi{\arabic{enumi}.}
\item
  \textbf{CMake} : Version \(\ge 3.20\).
\item
  \textbf{OpenMPI} : Version \(\ge 4.0\) (ou MPICH \(\ge 3.4\)). Fournit
  les exécutables \texttt{mpicxx} et \texttt{mpirun}.
\item
  \textbf{Python 3} : Version \(\ge 3.8\) avec les paquets scientifiques
  standards :

\begin{Shaded}
\begin{Highlighting}[]
\ExtensionTok{pip}\NormalTok{ install numpy matplotlib scipy}
\end{Highlighting}
\end{Shaded}
\end{enumerate}

\subsubsection{2.3. Procédure de Compilation pas-à-pas avec
CMake}\label{procuxe9dure-de-compilation-pas-uxe0-pas-avec-cmake}

Pour compiler COVOF 2D en mode optimisé pour la production :

\begin{Shaded}
\begin{Highlighting}[]
\CommentTok{\# 1. Se positionner à la racine du projet}
\BuiltInTok{cd}\NormalTok{ covof2d}

\CommentTok{\# 2. Créer le répertoire de compilation isolé}
\FunctionTok{mkdir} \AttributeTok{{-}p}\NormalTok{ build }\KeywordTok{\&\&} \BuiltInTok{cd}\NormalTok{ build}

\CommentTok{\# 3. Configurer le projet avec CMake en mode Release (optimisations {-}O3)}
\FunctionTok{cmake}\NormalTok{ .. }\AttributeTok{{-}DCMAKE\_BUILD\_TYPE}\OperatorTok{=}\NormalTok{Release }\AttributeTok{{-}DCMAKE\_CXX\_COMPILER}\OperatorTok{=}\NormalTok{mpicxx}

\CommentTok{\# 4. Compiler sur plusieurs cœurs (ex: 4 cœurs)}
\FunctionTok{make} \AttributeTok{{-}j4}

\CommentTok{\# 5. Vérifier la génération de l\textquotesingle{}exécutable}
\FunctionTok{ls} \AttributeTok{{-}lh}\NormalTok{ ../bin/covof2d}
\BuiltInTok{cd}\NormalTok{ ..}
\end{Highlighting}
\end{Shaded}

L'exécutable final est automatiquement placé dans le répertoire
\texttt{bin/covof2d}.

\subsubsection{2.4. Arborescence et Organisation des Fichiers du
Projet}\label{arborescence-et-organisation-des-fichiers-du-projet}

\begin{verbatim}
covof2d/
├── CMakeLists.txt              # Configuration de compilation CMake
├── bin/                        # Exécutable binaire final (covof2d)
├── include/covof/               # En-têtes C++ de la bibliothèque COVOF
│   ├── Config.hpp              # Structure de données de configuration
│   ├── Fields2D.hpp            # Allocation mémoire des variables scalaires et vectorielles
│   ├── InterfaceTracker.hpp    # Algorithmes VOF et modèle CSF capillaire
│   ├── LinearSolver.hpp        # Hiérarchie des solveurs de Krylov et préconditionneurs
│   ├── Mesh2D.hpp              # Grille cartésienne MAC décalée et métriques
│   ├── MpiComm.hpp             # Topologie cartésienne MPI et échanges de halos
│   ├── NavierStokesSolver.hpp  # Assemblage discret, Lagrangien Augmenté, Thermique
│   ├── Simulation.hpp          # Orchestration de l'avancement temporel
│   └── VTKWriter.hpp           # Génération des fichiers XML VTK pour ParaView
├── src/                        # Fichiers sources C++ implémentant les modules
│   ├── Fields2D.cpp
│   ├── InterfaceTracker.cpp
│   ├── LinearSolver.cpp
│   ├── Mesh2D.cpp
│   ├── MpiComm.cpp
│   ├── NavierStokesSolver.cpp
│   ├── Simulation.cpp
│   ├── VTKWriter.cpp
│   └── main.cpp                # Point d'entrée principal de l'exécutable
├── cases/                      # Fichiers de configuration JSON des benchmarks
├── docs/                       # Manuels complets de théorie, utilisateur, dev, validation
│   ├── MATHEMATICAL_THEORY.md  # Notice théorique et mathématique
│   ├── USER_GUIDE.md           # Ce manuel utilisateur
│   ├── DEVELOPER_GUIDE.md      # Guide d'architecture logicielle C++
│   ├── VALIDATION_GUIDE.md     # Cahier de validation et vérification Référence Académique CFD
│   └── figures/                # Figures haute définition des validations
└── scripts/                    # Scripts d'automatisation et de post-traitement
    ├── run_all_benchmarks.sh   # Lance la suite complète des 13 cas tests
    ├── verify_cases.py         # Analyse statistique et vérification des critères d'erreur
    └── generate_validation_figures.py # Générateur de figures PNG
\end{verbatim}

\begin{center}\rule{0.5\linewidth}{0.5pt}\end{center}

\subsection{3. Format Complet des Fichiers de Configuration
(JSON)}\label{format-complet-des-fichiers-de-configuration-json}

\subsubsection{3.1. Structure Générale d'un Cas
Test}\label{structure-guxe9nuxe9rale-dun-cas-test}

Chaque simulation est entièrement pilotée par un fichier texte structuré
selon le standard \textbf{JSON}. Un fichier type contient neuf blocs
fonctionnels :

\begin{Shaded}
\begin{Highlighting}[]
\FunctionTok{\{}
  \DataTypeTok{"case\_name"}\FunctionTok{:} \StringTok{"cavity\_re1000"}\FunctionTok{,}
  \DataTypeTok{"mesh"}\FunctionTok{:} \FunctionTok{\{} \ErrorTok{...} \FunctionTok{\},}
  \DataTypeTok{"physics"}\FunctionTok{:} \FunctionTok{\{} \ErrorTok{...} \FunctionTok{\},}
  \DataTypeTok{"thermal"}\FunctionTok{:} \FunctionTok{\{} \ErrorTok{...} \FunctionTok{\},}
  \DataTypeTok{"time"}\FunctionTok{:} \FunctionTok{\{} \ErrorTok{...} \FunctionTok{\},}
  \DataTypeTok{"boundary\_conditions"}\FunctionTok{:} \FunctionTok{\{} \ErrorTok{...} \FunctionTok{\},}
  \DataTypeTok{"immersed\_boundary"}\FunctionTok{:} \FunctionTok{\{} \ErrorTok{...} \FunctionTok{\},}
  \DataTypeTok{"initial\_condition"}\FunctionTok{:} \FunctionTok{\{} \ErrorTok{...} \FunctionTok{\},}
  \DataTypeTok{"solver"}\FunctionTok{:} \FunctionTok{\{} \ErrorTok{...} \FunctionTok{\},}
  \DataTypeTok{"output"}\FunctionTok{:} \FunctionTok{\{} \ErrorTok{...} \FunctionTok{\}}
\FunctionTok{\}}
\end{Highlighting}
\end{Shaded}

\subsubsection{3.2. Dictionnaire Exhaustif des Modules et
Mots-Clés}\label{dictionnaire-exhaustif-des-modules-et-mots-cluxe9s}

\paragraph{\texorpdfstring{3.2.1. Module \texttt{mesh} (Maillage
Cartésien)}{3.2.1. Module mesh (Maillage Cartésien)}}\label{module-mesh-maillage-cartuxe9sien}

Définit la géométrie du domaine rectangulaire et le nombre de cellules
de calcul : - \texttt{nx} (entier, \(>0\)) : Nombre de cellules de
contrôle dans la direction axiale \(X\). - \texttt{ny} (entier, \(>0\))
: Nombre de cellules de contrôle dans la direction transversale \(Y\). -
\texttt{lx} (réel, \(>0\)) : Longueur physique du domaine selon \(X\)
{[}m{]}. - \texttt{ly} (réel, \(>0\)) : Hauteur physique du domaine
selon \(Y\) {[}m{]}. - \texttt{origin\_x} (réel, optionnel, défaut:
\texttt{0.0}) : Abscisse du coin inférieur gauche {[}m{]}. -
\texttt{origin\_y} (réel, optionnel, défaut: \texttt{0.0}) : Ordonnée du
coin inférieur gauche {[}m{]}.

\paragraph{\texorpdfstring{3.2.2. Module \texttt{physics} (Propriétés et
Modèles
d'Écoulement)}{3.2.2. Module physics (Propriétés et Modèles d'Écoulement)}}\label{module-physics-propriuxe9tuxe9s-et-moduxe8les-duxe9coulement}

Contrôle les modèles physiques activés et les propriétés de fluide : -
\texttt{enable\_convection} (booléen) : Active (\texttt{true}) ou
désactive (\texttt{false}) les termes de transport convectif
\((\mathbf{u}\cdot\nabla)\mathbf{u}\). Doit être \texttt{false} pour les
écoulements de Stokes purs ou le problème de Szymanski. -
\texttt{convective\_scheme} (chaîne, défaut: \texttt{"UPWIND"}) : Schéma
numérique pour la convection : - \texttt{"UPWIND"} : Schéma décentré
amont d'ordre 1, inconditionnellement stable. - \texttt{"CENTERED"} :
Schéma centré d'ordre 2, non dissipatif. - \texttt{"TVD\_SUPERBEE"} :
Schéma d'ordre 2 avec limiteur de flux Superbee. - \texttt{gx} (réel,
défaut: \texttt{0.0}) : Accélération de la pesanteur selon \(X\)
{[}m/s²{]}. - \texttt{gy} (réel, défaut: \texttt{0.0}) : Accélération de
la pesanteur selon \(Y\) {[}m/s²{]} (ex: \texttt{-9.81} pour la gravité
terrestre). - \texttt{fluid1} (objet) : Propriétés thermodynamiques de
la phase principale : - \texttt{rho} (réel, \(>0\)) : Masse volumique
\([\text{kg/m}^3]\). - \texttt{mu} (réel, \(>0\)) : Viscosité dynamique
\([\text{Pa}\cdot\text{s}]\). - \texttt{enable\_two\_phase} (booléen,
défaut: \texttt{false}) : Active le modèle diphasique à deux fluides
immiscibles. - \texttt{fluid2} (objet, requis si
\texttt{enable\_two\_phase\ =\ true}) : Propriétés de la seconde phase
(\texttt{rho}, \texttt{mu}). - \texttt{enable\_surface\_tension}
(booléen, défaut: \texttt{false}) : Active le modèle de tension de
surface CSF de Brackbill. - \texttt{sigma} (réel, \(\ge 0\)) :
Coefficient de tension de surface \([\text{N/m}]\). -
\texttt{viscosity\_model} (chaîne, défaut: \texttt{"CONSTANT"}) : -
\texttt{"CONSTANT"} : Viscosité spatiale homogène par phase. -
\texttt{"EXPONENTIAL"} : Viscosité spatialement variable continue selon
la loi : \[\mu(y) = \mu_0 \exp\left( a \frac{(y-e)^2}{e^2} \right)\] -
\texttt{viscosity\_exp\_mu0}, \texttt{viscosity\_exp\_a},
\texttt{viscosity\_exp\_e} (réels) : Paramètres du modèle exponentiel.

\paragraph{\texorpdfstring{3.2.3. Module \texttt{thermal} (Couplage
Énergétique et
Boussinesq)}{3.2.3. Module thermal (Couplage Énergétique et Boussinesq)}}\label{module-thermal-couplage-uxe9nerguxe9tique-et-boussinesq}

Contrôle la résolution de l'équation de la chaleur et la force de
flottabilité : - \texttt{enable\_energy} (booléen, défaut:
\texttt{false}) : Active la résolution de l'équation d'énergie thermique
pour le champ de température \(T\). - \texttt{beta} (réel, défaut:
\texttt{0.0}) : Coefficient de dilatation volumique thermique
\([\text{K}^{-1}]\). - \texttt{t\_ref} (réel, défaut: \texttt{0.0}) :
Température de référence de Boussinesq \([\text{K}]\) ou
\([^{\circ}\text{C}]\). - \texttt{t\_hot} (réel) : Température imposée à
la paroi chaude \([\text{K}]\). - \texttt{t\_cold} (réel) : Température
imposée à la paroi froide \([\text{K}]\). - \texttt{k\_thermal} (réel,
\(>0\), défaut: \texttt{1.0}) : Conductivité thermique
\([\text{W/(m}\cdot\text{K)}]\). - \texttt{cp\_thermal} (réel, \(>0\),
défaut: \texttt{1.0}) : Capacité thermique isobare
\([\text{J/(kg}\cdot\text{K)}]\).

\paragraph{\texorpdfstring{3.2.4. Module \texttt{time} (Intégration
Temporelle)}{3.2.4. Module time (Intégration Temporelle)}}\label{module-time-intuxe9gration-temporelle}

\begin{itemize}
\tightlist
\item
  \texttt{dt} (réel, \(>0\)) : Pas de temps d'intégration
  \([\text{s}]\).
\item
  \texttt{t\_final} (réel, \(>0\)) : Temps physique maximal de la
  simulation \([\text{s}]\).
\item
  \texttt{adaptive\_dt} (booléen, défaut: \texttt{false}) : Active la
  régulation automatique du pas de temps selon la condition CFL.
\item
  \texttt{cfl\_max} (réel, défaut: \texttt{0.5}) : Nombre CFL limite
  admissible.
\end{itemize}

\paragraph{\texorpdfstring{3.2.5. Module \texttt{boundary\_conditions}
(Conditions aux
Limites)}{3.2.5. Module boundary\_conditions (Conditions aux Limites)}}\label{module-boundary_conditions-conditions-aux-limites}

Définit les conditions physiques sur les quatre frontières cardinales :
\texttt{left}, \texttt{right}, \texttt{bottom}, \texttt{top}. Chaque
frontière comporte : - \texttt{type} (chaîne) : - \texttt{"DIRICHLET"} :
Vitesse imposée à la frontière. Requièrt les champs \texttt{u} et
\texttt{v}. - \texttt{"FREE\_SLIP"} : Vitesse normale nulle, contrainte
de cisaillement nulle (symétrie). - \texttt{"NEUMANN"} : Flux de vitesse
normal nul (\(\partial \mathbf{u} / \partial n = \mathbf{0}\), sortie
libre). - \texttt{"PERIODIC"} : Périodicité spatiale liant les faces
opposées. - \texttt{thermal\_type} (chaîne, optionnel si la thermique
est activée) : - \texttt{"DIRICHLET"} : Température imposée égale à la
valeur \texttt{temperature}. - \texttt{"NEUMANN"} : Paroi adiabatique
isolée (\(\partial T / \partial n = 0\)).

\paragraph{\texorpdfstring{3.2.6. Module \texttt{immersed\_boundary}
(Pénalisation de
Brinkman)}{3.2.6. Module immersed\_boundary (Pénalisation de Brinkman)}}\label{module-immersed_boundary-puxe9nalisation-de-brinkman}

Permet d'insérer des obstacles solides rigides immobiles dans
l'écoulement : - \texttt{enabled} (booléen, défaut: \texttt{false}) :
Active la pénalisation. - \texttt{type} (chaîne) : - \texttt{"BOX"} :
Obstacle parallélépipédique (marche descendante, cavité en équerre).
Paramètres : \texttt{box\_x0}, \texttt{box\_y0}, \texttt{box\_x1},
\texttt{box\_y1}. - \texttt{"CYLINDER"} : Obstacle circulaire plein
(cylindre immergé). Paramètres : \texttt{cylinder\_xc},
\texttt{cylinder\_yc}, \texttt{cylinder\_r}. -
\texttt{permeability\_solid} (réel, défaut: \texttt{1.0e-40}) :
Perméabilité artificielle de Brinkman \(K_{solid}\) \([\text{m}^2]\).

\paragraph{\texorpdfstring{3.2.7. Module \texttt{initial\_condition}
(Champs
Initiaux)}{3.2.7. Module initial\_condition (Champs Initiaux)}}\label{module-initial_condition-champs-initiaux}

Définit la distribution initiale des variables à \(t=0\) : -
\texttt{type} (chaîne) : - \texttt{"UNIFORM"} : Vitesse initiale
uniforme \((\mathbf{u}_0, \mathbf{v}_0)\). - \texttt{"POISEUILLE"} :
Profil parabolique établi imposé comme condition initiale. -
\texttt{"DROPLET"} : Goutte circulaire de phase 1 dans une matrice de
phase 2 (\texttt{center\_x}, \texttt{center\_y}, \texttt{radius}). -
\texttt{"BUBBLE"} : Bulle circulaire de phase 2 dans un liquide lourd de
phase 1. - \texttt{"DAM\_BREAK"} : Colonne d'eau initiale définie par un
rectangle \([x_0, x_1] \times [y_0, y_1]\). - \texttt{"VORTEX\_SHEAR"} :
Disque fluide soumis à un cisaillement de déformation (Cas 3.19). -
\texttt{"RAYLEIGH\_TAYLOR"} : Stratification dense au-dessus d'un fluide
léger avec perturbation ondulatoire.

\paragraph{\texorpdfstring{3.2.8. Module \texttt{solver} (Solveurs
Algébriques et Lagrangien
Augmenté)}{3.2.8. Module solver (Solveurs Algébriques et Lagrangien Augmenté)}}\label{module-solver-solveurs-alguxe9briques-et-lagrangien-augmentuxe9}

\begin{itemize}
\tightlist
\item
  \texttt{type} (chaîne, défaut: \texttt{"AUGMENTED\_LAGRANGIAN"}) :
  Méthode de couplage vitesse-pression.
\item
  \texttt{linear\_solver} (chaîne, défaut: \texttt{"BICGSTAB2"}) :

  \begin{itemize}
  \tightlist
  \item
    \texttt{"BICGSTAB2"} : Bi-Conjugate Gradient Stabilized d'ordre 2
    (optimal).
  \item
    \texttt{"BICGSTAB"} : BiCGSTAB standard.
  \item
    \texttt{"CG"} : Gradient Conjugué (réservé aux matrices symétriques
    définies positives comme le Laplacien).
  \item
    \texttt{"GMRES"} : Generalized Minimal Residual avec redémarrages.
  \end{itemize}
\item
  \texttt{precond\_type} (chaîne, défaut: \texttt{"JACOBI"}) :

  \begin{itemize}
  \tightlist
  \item
    \texttt{"JACOBI"} : Préconditionnement diagonal (très performant et
    idéal en MPI).
  \item
    \texttt{"SSOR"} : Sur-relaxation successive symétrique.
  \item
    \texttt{"ILU0"} : Factorisation LU incomplète sans remplissage.
  \end{itemize}
\item
  \texttt{tol} (réel, défaut: \texttt{1.0e-5}) : Critère d'arrêt sur le
  résidu relatif \(\|\mathbf{r}\| / \|\mathbf{b}\|\).
\item
  \texttt{max\_iter} (entier, défaut: \texttt{300}) : Nombre maximal
  d'itérations de Krylov autorisées par pas.
\item
  \texttt{r\_al} (réel, défaut: \texttt{500.0}) : Paramètre de pénalité
  de divergence \(r\) \([\text{kg/(m}^3\cdot\text{s)}]\).
\item
  \texttt{uzawa\_iter} (entier, défaut: \texttt{1}) : Nombre
  d'itérations d'Uzawa par pas de temps.
\end{itemize}

\paragraph{\texorpdfstring{3.2.9. Module \texttt{output} (Sauvegardes et
Visualisation)}{3.2.9. Module output (Sauvegardes et Visualisation)}}\label{module-output-sauvegardes-et-visualisation}

\begin{itemize}
\tightlist
\item
  \texttt{directory} (chaîne) : Nom du dossier où sont enregistrés les
  fichiers de sortie.
\item
  \texttt{vtk\_interval} (entier) : Intervalle d'écriture des fichiers
  ParaView \texttt{.pvtr} / \texttt{.vtr} (en nombre de pas).
\item
  \texttt{csv\_interval} (entier) : Intervalle d'écriture dans
  l'historique scalaire \texttt{.csv}.
\item
  \texttt{console\_interval} (entier) : Intervalle de rafraîchissement
  des informations dans le terminal.
\end{itemize}

\begin{center}\rule{0.5\linewidth}{0.5pt}\end{center}

\subsection{4. Tutoriel pas-à-pas : Construire un Cas de
Simulation}\label{tutoriel-pas-uxe0-pas-construire-un-cas-de-simulation}

\subsubsection{\texorpdfstring{4.1. Tutoriel 1 : Écoulement Monophasique
dans une Cavité Entraînée
(\(Re=1000\))}{4.1. Tutoriel 1 : Écoulement Monophasique dans une Cavité Entraînée (Re=1000)}}\label{tutoriel-1-uxe9coulement-monophasique-dans-une-cavituxe9-entrauxeenuxe9e-re1000}

Dans ce cas classique, un fluide visqueux est enfermé dans une cavité
carrée de \(1\text{ m} \times 1\text{ m}\). Le couvercle supérieur se
déplace à vitesse horizontale \(u_{lid} = 1.0\text{ m/s}\). La viscosité
est ajustée à \(\mu = 10^{-3}\text{ Pa}\cdot\text{s}\) avec une densité
\(\rho = 1\text{ kg/m}^3\) pour obtenir
\(Re = \frac{\rho U L}{\mu} = 1000\).

Créez le fichier \texttt{cases/mon\_cas\_cavite.json} :

\begin{Shaded}
\begin{Highlighting}[]
\FunctionTok{\{}
  \DataTypeTok{"case\_name"}\FunctionTok{:} \StringTok{"cavite\_tutoriel"}\FunctionTok{,}
  \DataTypeTok{"mesh"}\FunctionTok{:} \FunctionTok{\{}
    \DataTypeTok{"nx"}\FunctionTok{:} \DecValTok{128}\FunctionTok{,}
    \DataTypeTok{"ny"}\FunctionTok{:} \DecValTok{128}\FunctionTok{,}
    \DataTypeTok{"lx"}\FunctionTok{:} \FloatTok{1.0}\FunctionTok{,}
    \DataTypeTok{"ly"}\FunctionTok{:} \FloatTok{1.0}
  \FunctionTok{\},}
  \DataTypeTok{"physics"}\FunctionTok{:} \FunctionTok{\{}
    \DataTypeTok{"enable\_convection"}\FunctionTok{:} \KeywordTok{true}\FunctionTok{,}
    \DataTypeTok{"enable\_two\_phase"}\FunctionTok{:} \KeywordTok{false}\FunctionTok{,}
    \DataTypeTok{"enable\_surface\_tension"}\FunctionTok{:} \KeywordTok{false}\FunctionTok{,}
    \DataTypeTok{"gx"}\FunctionTok{:} \FloatTok{0.0}\FunctionTok{,}
    \DataTypeTok{"gy"}\FunctionTok{:} \FloatTok{0.0}\FunctionTok{,}
    \DataTypeTok{"fluid1"}\FunctionTok{:} \FunctionTok{\{}
      \DataTypeTok{"rho"}\FunctionTok{:} \FloatTok{1.0}\FunctionTok{,}
      \DataTypeTok{"mu"}\FunctionTok{:} \FloatTok{0.001}
    \FunctionTok{\}}
  \FunctionTok{\},}
  \DataTypeTok{"time"}\FunctionTok{:} \FunctionTok{\{}
    \DataTypeTok{"dt"}\FunctionTok{:} \FloatTok{0.05}\FunctionTok{,}
    \DataTypeTok{"t\_final"}\FunctionTok{:} \FloatTok{40.0}\FunctionTok{,}
    \DataTypeTok{"adaptive\_dt"}\FunctionTok{:} \KeywordTok{false}
  \FunctionTok{\},}
  \DataTypeTok{"boundary\_conditions"}\FunctionTok{:} \FunctionTok{\{}
    \DataTypeTok{"left"}\FunctionTok{:}   \FunctionTok{\{} \DataTypeTok{"type"}\FunctionTok{:} \StringTok{"DIRICHLET"}\FunctionTok{,} \DataTypeTok{"u"}\FunctionTok{:} \FloatTok{0.0}\FunctionTok{,} \DataTypeTok{"v"}\FunctionTok{:} \FloatTok{0.0} \FunctionTok{\},}
    \DataTypeTok{"right"}\FunctionTok{:}  \FunctionTok{\{} \DataTypeTok{"type"}\FunctionTok{:} \StringTok{"DIRICHLET"}\FunctionTok{,} \DataTypeTok{"u"}\FunctionTok{:} \FloatTok{0.0}\FunctionTok{,} \DataTypeTok{"v"}\FunctionTok{:} \FloatTok{0.0} \FunctionTok{\},}
    \DataTypeTok{"bottom"}\FunctionTok{:} \FunctionTok{\{} \DataTypeTok{"type"}\FunctionTok{:} \StringTok{"DIRICHLET"}\FunctionTok{,} \DataTypeTok{"u"}\FunctionTok{:} \FloatTok{0.0}\FunctionTok{,} \DataTypeTok{"v"}\FunctionTok{:} \FloatTok{0.0} \FunctionTok{\},}
    \DataTypeTok{"top"}\FunctionTok{:}    \FunctionTok{\{} \DataTypeTok{"type"}\FunctionTok{:} \StringTok{"DIRICHLET"}\FunctionTok{,} \DataTypeTok{"u"}\FunctionTok{:} \FloatTok{1.0}\FunctionTok{,} \DataTypeTok{"v"}\FunctionTok{:} \FloatTok{0.0} \FunctionTok{\}}
  \FunctionTok{\},}
  \DataTypeTok{"solver"}\FunctionTok{:} \FunctionTok{\{}
    \DataTypeTok{"type"}\FunctionTok{:} \StringTok{"AUGMENTED\_LAGRANGIAN"}\FunctionTok{,}
    \DataTypeTok{"linear\_solver"}\FunctionTok{:} \StringTok{"BICGSTAB2"}\FunctionTok{,}
    \DataTypeTok{"precond\_type"}\FunctionTok{:} \StringTok{"JACOBI"}\FunctionTok{,}
    \DataTypeTok{"tol"}\FunctionTok{:} \FloatTok{1.0e{-}5}\FunctionTok{,}
    \DataTypeTok{"max\_iter"}\FunctionTok{:} \DecValTok{300}\FunctionTok{,}
    \DataTypeTok{"r\_al"}\FunctionTok{:} \FloatTok{100.0}\FunctionTok{,}
    \DataTypeTok{"uzawa\_iter"}\FunctionTok{:} \DecValTok{1}
  \FunctionTok{\},}
  \DataTypeTok{"output"}\FunctionTok{:} \FunctionTok{\{}
    \DataTypeTok{"directory"}\FunctionTok{:} \StringTok{"output\_tutoriel\_cavite"}\FunctionTok{,}
    \DataTypeTok{"vtk\_interval"}\FunctionTok{:} \DecValTok{100}\FunctionTok{,}
    \DataTypeTok{"console\_interval"}\FunctionTok{:} \DecValTok{20}
  \FunctionTok{\}}
\FunctionTok{\}}
\end{Highlighting}
\end{Shaded}

Lancez le calcul sur 4 cœurs MPI :

\begin{Shaded}
\begin{Highlighting}[]
\ExtensionTok{mpirun} \AttributeTok{{-}np}\NormalTok{ 4 bin/covof2d cases/mon\_cas\_cavite.json}
\end{Highlighting}
\end{Shaded}

\begin{center}\rule{0.5\linewidth}{0.5pt}\end{center}

\subsection{5. Exécution du Code en Séquentiel et Parallèle
(OpenMPI)}\label{exuxe9cution-du-code-en-suxe9quentiel-et-paralluxe8le-openmpi}

\subsubsection{5.1. Lancement Local sur Poste de
Travail}\label{lancement-local-sur-poste-de-travail}

Pour exécuter un cas en séquentiel :

\begin{Shaded}
\begin{Highlighting}[]
\ExtensionTok{bin/covof2d}\NormalTok{ cases/mon\_cas.json}
\end{Highlighting}
\end{Shaded}

Pour exécuter en parallèle sur plusieurs cœurs physiques de votre
machine :

\begin{Shaded}
\begin{Highlighting}[]
\ExtensionTok{mpirun} \AttributeTok{{-}np}\NormalTok{ 2 bin/covof2d cases/mon\_cas.json}
\ExtensionTok{mpirun} \AttributeTok{{-}np}\NormalTok{ 8 bin/covof2d cases/mon\_cas.json}
\end{Highlighting}
\end{Shaded}

\subsubsection{5.2. Lancement sur Supercalculateur et Gestionnaire de
Tâches
SLURM}\label{lancement-sur-supercalculateur-et-gestionnaire-de-tuxe2ches-slurm}

Sur un cluster de calcul administré par SLURM, créez un script de
soumission \texttt{job\_covof2d.sh} :

\begin{Shaded}
\begin{Highlighting}[]
\CommentTok{\#!/bin/bash}
\CommentTok{\#SBATCH {-}{-}job{-}name=covof2d\_calc}
\CommentTok{\#SBATCH {-}{-}output=covof2d\_\%j.log}
\CommentTok{\#SBATCH {-}{-}error=covof2d\_\%j.err}
\CommentTok{\#SBATCH {-}{-}nodes=2}
\CommentTok{\#SBATCH {-}{-}ntasks{-}per{-}node=32}
\CommentTok{\#SBATCH {-}{-}time=04:00:00}
\CommentTok{\#SBATCH {-}{-}partition=standard}

\ExtensionTok{module}\NormalTok{ purge}
\ExtensionTok{module}\NormalTok{ load gcc/11.2.0 openmpi/4.1.2 cmake/3.22.1}

\BuiltInTok{export} \VariableTok{OMP\_NUM\_THREADS}\OperatorTok{=}\NormalTok{1}

\ExtensionTok{srun}\NormalTok{ bin/covof2d cases/cavity\_re1000.json}
\end{Highlighting}
\end{Shaded}

Soumettez la tâche :

\begin{Shaded}
\begin{Highlighting}[]
\ExtensionTok{sbatch}\NormalTok{ job\_covof2d.sh}
\end{Highlighting}
\end{Shaded}

\subsubsection{5.3. Décomposition de Domaine MPI et Équilibrage de
Charge}\label{duxe9composition-de-domaine-mpi-et-uxe9quilibrage-de-charge}

COVOF 2D divise automatiquement la grille cartésienne globale
\(N_x \times N_y\) en une grille de sous-domaines
\(P_x \times P_y = NP\) de façon à minimiser le rapport
périmètre/surface (minimisation des volumes de communication
inter-processus).

Chaque processus MPI gère son propre sous-domaine local bordé par une
rangée de cellules fantômes (\emph{halos}) qui sont synchronisées de
façon asynchrone par \texttt{MPI\_Isend} et \texttt{MPI\_Irecv} à chaque
étape de produit matrice-vecteur et à chaque itération temporelle.

\begin{center}\rule{0.5\linewidth}{0.5pt}\end{center}

\subsection{6. Suivi d'Exécution et Diagnostic des Calculs en Temps
Réel}\label{suivi-dexuxe9cution-et-diagnostic-des-calculs-en-temps-ruxe9el}

\subsubsection{6.1. Interprétation du Tableau de Bord
Console}\label{interpruxe9tation-du-tableau-de-bord-console}

Lors de l'exécution, COVOF 2D affiche périodiquement un tableau
synthétique :

\begin{verbatim}
Pas     Temps [s]     dt [s]      Iter Krylov Vitesse Max     Div Max         Nu Moyen
-----------------------------------------------------------------------------------------
50      2.50000       5.000e-02   78          9.4120e-01      3.1205e-04      2.7261
100     5.00000       5.000e-02   64          9.4350e-01      1.8410e-04      2.7261
\end{verbatim}

\begin{itemize}
\tightlist
\item
  \textbf{Pas} : Numéro du pas de temps courant.
\item
  \textbf{Temps {[}s{]}} : Temps physique accumulé depuis le début de la
  simulation.
\item
  \textbf{dt {[}s{]}} : Pas de temps appliqué.
\item
  \textbf{Iter Krylov} : Nombre d'itérations nécessaires au solveur
  BiCGSTAB(2) pour atteindre la tolérance résiduelle.
\item
  \textbf{Vitesse Max {[}m/s{]}} : Norme infinie
  \(\|\mathbf{u}\|_\infty\) du champ de vitesse.
\item
  \textbf{Div Max {[}s⁻¹{]}} : Norme infinie de la divergence discrète
  \(\|\nabla \cdot \mathbf{u}\|_\infty\). Elle doit rester inférieure à
  \(10^{-3}\text{ s}^{-1}\) pour garantir la stricte incompressibilité
  physique.
\item
  \textbf{Nu Moyen} : Nombre de Nusselt moyen à la paroi chaude (affiché
  si le module thermique est activé).
\end{itemize}

\subsubsection{6.2. Détection Précoce des Instabilités et
Divergences}\label{duxe9tection-pruxe9coce-des-instabilituxe9s-et-divergences}

\begin{enumerate}
\def\labelenumi{\arabic{enumi}.}
\tightlist
\item
  \textbf{Explosion de la Vitesse Max} : Si la vitesse dépasse
  brutalement plusieurs fois la vitesse caractéristique du problème, la
  condition CFL a été enfreinte ou le terme convectif centré est devenu
  instable.
\item
  \textbf{Stagnation des itérations Krylov à \texttt{max\_iter}} :
  Indique un mauvais conditionnement de la matrice, généralement résolu
  en augmentant le paramètre de pénalité \(r_{al}\) ou en diminuant
  \(\Delta t\).
\item
  \textbf{Div Max anormalement élevée (\(> 0.1\))} : Le paramètre de
  pénalité \(r_{al}\) est trop faible, ou un plus grand nombre
  d'itérations d'Uzawa (\texttt{uzawa\_iter\ =\ 2}) est requis.
\end{enumerate}

\begin{center}\rule{0.5\linewidth}{0.5pt}\end{center}

\subsection{7. Post-Traitement et Analyse Graphique avec
ParaView}\label{post-traitement-et-analyse-graphique-avec-paraview}

\subsubsection{\texorpdfstring{7.1. Chargement des Fichiers
\texttt{.pvtr} et
\texttt{.vtr}}{7.1. Chargement des Fichiers .pvtr et .vtr}}\label{chargement-des-fichiers-.pvtr-et-.vtr}

\begin{enumerate}
\def\labelenumi{\arabic{enumi}.}
\tightlist
\item
  Lancez ParaView (\texttt{paraview}).
\item
  Cliquez sur \textbf{File -\textgreater{} Open\ldots{}} et naviguez
  vers le dossier de résultats (ex: \texttt{output\_cavity\_re1000}).
\item
  Sélectionnez le groupe de fichiers maîtres :
  \texttt{cavity\_re1000\_..pvtr}.
\item
  Cliquez sur \textbf{Apply} dans le panneau \emph{Properties}.
\end{enumerate}

\subsubsection{7.2. Extraction des Lignes de Courant et Profils de
Vitesse}\label{extraction-des-lignes-de-courant-et-profils-de-vitesse}

\begin{itemize}
\tightlist
\item
  \textbf{Lignes de courant} :

  \begin{enumerate}
  \def\labelenumi{\arabic{enumi}.}
  \tightlist
  \item
    Sélectionnez votre source dans l'arbre d'objets.
  \item
    Appliquez le filtre \textbf{Stream Tracer}.
  \item
    Choisissez le vecteur \texttt{Velocity}, positionnez la ligne source
    de graines (Seed Type: High Resolution Line Source) et cliquez sur
    \textbf{Apply}.
  \end{enumerate}
\item
  \textbf{Profils axiaux 1D} :

  \begin{enumerate}
  \def\labelenumi{\arabic{enumi}.}
  \tightlist
  \item
    Appliquez le filtre \textbf{Plot Over Line}.
  \item
    Définissez les coordonnées de début \((0.5, 0.0)\) et de fin
    \((0.5, 1.0)\).
  \item
    Cliquez sur \textbf{Apply} pour visualiser la courbe \(u(y)\) à
    mi-cavité.
  \end{enumerate}
\end{itemize}

\begin{center}\rule{0.5\linewidth}{0.5pt}\end{center}

\subsection{8. Automatisation des Benchmarks et Scripts de
Vérification}\label{automatisation-des-benchmarks-et-scripts-de-vuxe9rification}

COVOF 2D est livré avec une suite d'outils Python et Bash permettant de
reproduire l'ensemble des cas tests académiques en un clic :

\subsubsection{\texorpdfstring{8.1. Le Script d'Orchestration Globale
\texttt{run\_all\_benchmarks.sh}}{8.1. Le Script d'Orchestration Globale run\_all\_benchmarks.sh}}\label{le-script-dorchestration-globale-run_all_benchmarks.sh}

\begin{Shaded}
\begin{Highlighting}[]
\CommentTok{\# Exécution complète de toute la suite avec OpenMPI sur 2 processus :}
\FunctionTok{bash}\NormalTok{ scripts/run\_all\_benchmarks.sh 2}
\end{Highlighting}
\end{Shaded}

Ce script compile le code, vide les anciens répertoires de sortie,
exécute successivement les 13 cas tests, lance le script de vérification
d'erreur et génère l'ensemble des figures.

\subsubsection{\texorpdfstring{8.2. Vérification Statistique Automatisée
\texttt{verify\_cases.py}}{8.2. Vérification Statistique Automatisée verify\_cases.py}}\label{vuxe9rification-statistique-automatisuxe9e-verify_cases.py}

\begin{Shaded}
\begin{Highlighting}[]
\ExtensionTok{python3}\NormalTok{ scripts/verify\_cases.py}
\end{Highlighting}
\end{Shaded}

Lit les résultats numériques dans les dossiers \texttt{output\_*},
calcule les normes d'erreur \(L_1, L_2, L_\infty\), les confronte aux
valeurs de référence du Cahier de Validation et Vérification CFD 2D et
produit un tableau récapitulatif \texttt{validation\_summary.csv} et
\texttt{validation\_summary.md}.

\subsubsection{\texorpdfstring{8.3. Production Automatique des Figures
de Rapport
\texttt{generate\_validation\_figures.py}}{8.3. Production Automatique des Figures de Rapport generate\_validation\_figures.py}}\label{production-automatique-des-figures-de-rapport-generate_validation_figures.py}

\begin{Shaded}
\begin{Highlighting}[]
\ExtensionTok{python3}\NormalTok{ scripts/generate\_validation\_figures.py}
\end{Highlighting}
\end{Shaded}

Reconstitue les sous-domaines MPI partagés, extrait les champs scalaires
et vectoriels et trace les 14 figures de référence enregistrées dans
\texttt{docs/figures/}.

\begin{center}\rule{0.5\linewidth}{0.5pt}\end{center}

\subsection{9. Foire Aux Questions (FAQ) et Astuces
Pratiques}\label{foire-aux-questions-faq-et-astuces-pratiques}

\textbf{Q : Comment choisir au mieux le paramètre de pénalité \(r_{al}\)
?}\\
\emph{R :} En règle générale, choisissez
\(r_{al} \approx 100 \times \frac{\mu}{\rho}\) pour des régimes
laminaires standards, ou
\(r_{al} \approx \frac{\rho \Delta x^2}{\Delta t}\) pour des calculs
transitoires fins. Une valeur trop petite nuit à l'incompressibilité
(\(\text{div } \mathbf{u} > 0\)), tandis qu'une valeur démesurée
(\(r_{al} > 10^7\)) dégrade le conditionnement de la matrice.

\textbf{Q : Pourquoi privilégier BiCGSTAB(2) plutôt que GMRES ?}\\
\emph{R :} BiCGSTAB(2) ne requiert le stockage que de quelques vecteurs
de travail (empreinte mémoire minime et fixe), tout en éliminant les
décrochages de convergence propres au BiCGSTAB simple. GMRES nécessite
de stocker la base de Krylov complète, ce qui augmente fortement la
consommation de RAM pour de grands maillages.

\textbf{Q : Comment modéliser un profil de vitesse pulsé ou
instationnaire à l'entrée ?}\\
\emph{R :} Consultez le
\href{file:///Users/elhadjiibrahimathiam/Documents/PROJET_ANTIGRAVITY/covof2d/docs/DEVELOPER_GUIDE.md}{Manuel
Développeur} pour surcharger la méthode
\texttt{apply\_velocity\_boundary\_conditions} dans
\texttt{src/MpiComm.cpp}.

\begin{center}\rule{0.5\linewidth}{0.5pt}\end{center}

\subsection{10. Guide Détaillé de Dépannage et Traitement des
Erreurs}\label{guide-duxe9tailluxe9-de-duxe9pannage-et-traitement-des-erreurs}

Lors de calculs complexes sur maillages fins ou à haut nombre de
Reynolds / Rayleigh, des difficultés numériques peuvent subvenir. Ce
tableau répertorie les 12 cas d'incidents les plus fréquents et leur
protocole de résolution :

{\def\LTcaptype{none} % do not increment counter
\begin{longtable}[]{@{}
  >{\raggedright\arraybackslash}p{(\linewidth - 6\tabcolsep) * \real{0.2500}}
  >{\raggedright\arraybackslash}p{(\linewidth - 6\tabcolsep) * \real{0.2500}}
  >{\raggedright\arraybackslash}p{(\linewidth - 6\tabcolsep) * \real{0.2500}}
  >{\raggedright\arraybackslash}p{(\linewidth - 6\tabcolsep) * \real{0.2500}}@{}}
\toprule\noalign{}
\begin{minipage}[b]{\linewidth}\raggedright
Symptôme / Message d'Erreur
\end{minipage} & \begin{minipage}[b]{\linewidth}\raggedright
Cause Numérique ou Physique
\end{minipage} & \begin{minipage}[b]{\linewidth}\raggedright
Diagnostic Immédiat
\end{minipage} & \begin{minipage}[b]{\linewidth}\raggedright
Action de Remédiation Recommandée
\end{minipage} \\
\midrule\noalign{}
\endhead
\bottomrule\noalign{}
\endlastfoot
\texttt{Krylov\ solver\ failed\ to\ converge\ (max\_iter\ reached)} &
Conditionnement dégradé de la matrice & Le résidu
\(\|\mathbf{r}\|/\|\mathbf{b}\|\) stagne au-dessus de \texttt{tol}. & 1.
Réduire le pas de temps \texttt{dt} d'un facteur 2.2. Augmenter le
paramètre de pénalité \texttt{r\_al} (ex: passer de 100 à 500).3.
Augmenter \texttt{max\_iter} (passer de 300 à 600).4. Basculer de
\texttt{BICGSTAB} à \texttt{BICGSTAB2} ou \texttt{GMRES}. \\
\texttt{Floating\ Point\ Exception\ (NaN\ or\ Inf\ detected)} &
Dépassement brutal de la limite de stabilité CFL & Vitesse ou gradient
explose vers l'infini. & 1. Vérifier la vitesse imposée aux
frontières.2. Activer \texttt{adaptive\_dt\ =\ true} avec
\texttt{cfl\_max\ =\ 0.3}.3. Basculer le schéma convectif de
\texttt{CENTERED} vers \texttt{UPWIND} ou \texttt{TVD\_SUPERBEE}. \\
\texttt{Divergence\ max\ reste\ élevée\ (\textgreater{}\ 0.01\ s\^{}-1)}
& Incompressibilité insuffisante & La contrainte
\(\nabla \cdot \mathbf{u} = 0\) n'est pas satisfaite. & 1. Augmenter
\texttt{uzawa\_iter} à 2 ou 3.2. Augmenter \texttt{r\_al} d'un facteur
10.3. Resserrer la tolérance du solveur \texttt{tol} à \(10^{-6}\). \\
\texttt{Oscillations\ parasites\ de\ l\textquotesingle{}interface\ VOF}
& Violation de la restriction capillaire de Brackbill & Ondes de haute
fréquence déforment la bulle/goutte. & 1. Vérifier que
\(\Delta t \le \sqrt{\frac{(\rho_1+\rho_2)\Delta x^3}{4\pi\sigma}}\).2.
Réduire \texttt{dt} d'un facteur 2.3. Vérifier que la tension de surface
\texttt{sigma} est réaliste. \\
\texttt{Perte\ de\ masse\ excessive\ de\ la\ bulle\ (\textgreater{}\ 3\%)}
& Diffusion numérique ou schéma convectif d'ordre 1 & Le volume de la
phase VOF diminue au cours du temps. & 1. Raffiner le maillage spatial
dans la zone interfaciale.2. Utiliser le limiteur \texttt{TVD\_SUPERBEE}
pour préserver la netteté du front. \\
\texttt{MPI\ Abort\ /\ Deadlock\ à\ l\textquotesingle{}échange\ de\ halos}
& Blocage dans les communications point-à-point & Un processus attend un
message non émis. & 1. Vérifier la parité de la grille cartésienne
MPI.2. S'assurer que le nombre de cellules locales \(N_x^{loc} \ge 2\)
par processus. \\
\texttt{Vitesse\ non\ nulle\ dans\ un\ obstacle\ solide\ immergé} &
Perméabilité de Brinkman \(K_{solid}\) trop élevée & Le fluide pénètre à
l'intérieur du corps solide. & 1. Vérifier que
\texttt{permeability\_solid} est fixée à \(10^{-40}\text{ m}^2\).2.
Vérifier la fonction d'appartenance géométrique (\texttt{BOX} ou
\texttt{CYLINDER}). \\
\texttt{Profil\ de\ température\ instable\ ou\ T\ \textless{}\ T\_cold}
& Dispersion thermique ou condition limite incohérente & Valeurs hors
bornes thermodynamiques. & 1. S'assurer que les conditions limites
thermiques respectent le principe du maximum.2. Utiliser une
discrétisation implicite pour l'énergie thermique. \\
\texttt{Segmentation\ Fault\ au\ démarrage} & Dépassement de la mémoire
RAM disponible & La grille globale \(N_x \times N_y\) excède la capacité
mémoire. & 1. Augmenter le nombre de processus MPI pour répartir
l'empreinte mémoire.2. Réduire la taille de la grille
\(N_x \times N_y\). \\
\texttt{Nombre\ de\ Nusselt\ oscillant\ ou\ divergent} & Maillage
insuffisant dans la couche limite thermique & Le gradient pariétal
\(\partial T / \partial x\) est mal résolu. & 1. Raffiner la grille
cartésienne en \(X\) pour disposer d'au moins 8 mailles dans la couche
limite thermique \(\delta_T \sim L Ra^{-1/4}\). \\
\texttt{Courants\ parasites\ importants\ autour\ d\textquotesingle{}une\ goutte\ statique}
& Déséquilibre entre gradient de pression et force CSF & Des vortex
artificiels apparaissent autour de la goutte au repos. & 1. Vérifier que
la formulation Balanced-Force est active.2. S'assurer que la grille est
strictement carrée (\(\Delta x = \Delta y\)). \\
\texttt{Erreur\ de\ lecture\ JSON\ (SyntaxError\ or\ Missing\ Key)} &
Format JSON corrompu ou virgule orpheline & Le parseur rejette le
fichier de configuration. & 1. Valider le fichier JSON avec l'utilitaire
\texttt{python3\ -m\ json.tool\ cases/mon\_cas.json}. \\
\end{longtable}
}

\begin{center}\rule{0.5\linewidth}{0.5pt}\end{center}

\subsection{11. Annexe Exhaustive : Fiches de Paramétrage des 13 Cas
Tests}\label{annexe-exhaustive-fiches-de-paramuxe9trage-des-13-cas-tests}

Chaque cas de la suite de validation Référence Académique CFD dispose
d'un fichier JSON complet et prêt à l'emploi dans le répertoire
\texttt{cases/} :

\subsubsection{\texorpdfstring{11.1. Cas 3.1 :
\texttt{cases/laplacian.json} (Laplacien Cartésien
\(\nabla^2 \psi = 1\))}{11.1. Cas 3.1 : cases/laplacian.json (Laplacien Cartésien \textbackslash nabla\^{}2 \textbackslash psi = 1)}}\label{cas-3.1-caseslaplacian.json-laplacien-cartuxe9sien-nabla2-psi-1}

\begin{Shaded}
\begin{Highlighting}[]
\FunctionTok{\{}
  \DataTypeTok{"case\_name"}\FunctionTok{:} \StringTok{"laplacian"}\FunctionTok{,}
  \DataTypeTok{"mesh"}\FunctionTok{:} \FunctionTok{\{} \DataTypeTok{"nx"}\FunctionTok{:} \DecValTok{64}\FunctionTok{,} \DataTypeTok{"ny"}\FunctionTok{:} \DecValTok{32}\FunctionTok{,} \DataTypeTok{"lx"}\FunctionTok{:} \FloatTok{2.0}\FunctionTok{,} \DataTypeTok{"ly"}\FunctionTok{:} \FloatTok{1.0} \FunctionTok{\},}
  \DataTypeTok{"physics"}\FunctionTok{:} \FunctionTok{\{}
    \DataTypeTok{"enable\_convection"}\FunctionTok{:} \KeywordTok{false}\FunctionTok{,}
    \DataTypeTok{"enable\_two\_phase"}\FunctionTok{:} \KeywordTok{false}\FunctionTok{,}
    \DataTypeTok{"enable\_surface\_tension"}\FunctionTok{:} \KeywordTok{false}\FunctionTok{,}
    \DataTypeTok{"gx"}\FunctionTok{:} \FloatTok{0.0}\FunctionTok{,} \DataTypeTok{"gy"}\FunctionTok{:} \FloatTok{0.0}\FunctionTok{,}
    \DataTypeTok{"fluid1"}\FunctionTok{:} \FunctionTok{\{} \DataTypeTok{"rho"}\FunctionTok{:} \FloatTok{1.0}\FunctionTok{,} \DataTypeTok{"mu"}\FunctionTok{:} \FloatTok{1.0} \FunctionTok{\}}
  \FunctionTok{\},}
  \DataTypeTok{"time"}\FunctionTok{:} \FunctionTok{\{} \DataTypeTok{"dt"}\FunctionTok{:} \FloatTok{1.0}\FunctionTok{,} \DataTypeTok{"t\_final"}\FunctionTok{:} \FloatTok{1.0}\FunctionTok{,} \DataTypeTok{"adaptive\_dt"}\FunctionTok{:} \KeywordTok{false} \FunctionTok{\},}
  \DataTypeTok{"boundary\_conditions"}\FunctionTok{:} \FunctionTok{\{}
    \DataTypeTok{"left"}\FunctionTok{:}   \FunctionTok{\{} \DataTypeTok{"type"}\FunctionTok{:} \StringTok{"DIRICHLET"}\FunctionTok{,} \DataTypeTok{"u"}\FunctionTok{:} \FloatTok{0.0}\FunctionTok{,} \DataTypeTok{"v"}\FunctionTok{:} \FloatTok{0.0} \FunctionTok{\},}
    \DataTypeTok{"right"}\FunctionTok{:}  \FunctionTok{\{} \DataTypeTok{"type"}\FunctionTok{:} \StringTok{"DIRICHLET"}\FunctionTok{,} \DataTypeTok{"u"}\FunctionTok{:} \FloatTok{0.0}\FunctionTok{,} \DataTypeTok{"v"}\FunctionTok{:} \FloatTok{0.0} \FunctionTok{\},}
    \DataTypeTok{"bottom"}\FunctionTok{:} \FunctionTok{\{} \DataTypeTok{"type"}\FunctionTok{:} \StringTok{"DIRICHLET"}\FunctionTok{,} \DataTypeTok{"u"}\FunctionTok{:} \FloatTok{0.0}\FunctionTok{,} \DataTypeTok{"v"}\FunctionTok{:} \FloatTok{0.0} \FunctionTok{\},}
    \DataTypeTok{"top"}\FunctionTok{:}    \FunctionTok{\{} \DataTypeTok{"type"}\FunctionTok{:} \StringTok{"DIRICHLET"}\FunctionTok{,} \DataTypeTok{"u"}\FunctionTok{:} \FloatTok{0.0}\FunctionTok{,} \DataTypeTok{"v"}\FunctionTok{:} \FloatTok{0.0} \FunctionTok{\}}
  \FunctionTok{\},}
  \DataTypeTok{"initial\_condition"}\FunctionTok{:} \FunctionTok{\{} \DataTypeTok{"type"}\FunctionTok{:} \StringTok{"LAPLACIAN\_TEST"} \FunctionTok{\},}
  \DataTypeTok{"solver"}\FunctionTok{:} \FunctionTok{\{}
    \DataTypeTok{"type"}\FunctionTok{:} \StringTok{"AUGMENTED\_LAGRANGIAN"}\FunctionTok{,}
    \DataTypeTok{"linear\_solver"}\FunctionTok{:} \StringTok{"CG"}\FunctionTok{,}
    \DataTypeTok{"precond\_type"}\FunctionTok{:} \StringTok{"JACOBI"}\FunctionTok{,}
    \DataTypeTok{"tol"}\FunctionTok{:} \FloatTok{1.0e{-}12}\FunctionTok{,} \DataTypeTok{"max\_iter"}\FunctionTok{:} \DecValTok{500}\FunctionTok{,} \DataTypeTok{"r\_al"}\FunctionTok{:} \FloatTok{0.0}\FunctionTok{,} \DataTypeTok{"uzawa\_iter"}\FunctionTok{:} \DecValTok{1}
  \FunctionTok{\},}
  \DataTypeTok{"output"}\FunctionTok{:} \FunctionTok{\{} \DataTypeTok{"directory"}\FunctionTok{:} \StringTok{"output\_laplacian"}\FunctionTok{,} \DataTypeTok{"vtk\_interval"}\FunctionTok{:} \DecValTok{1}\FunctionTok{,} \DataTypeTok{"console\_interval"}\FunctionTok{:} \DecValTok{1} \FunctionTok{\}}
\FunctionTok{\}}
\end{Highlighting}
\end{Shaded}

\subsubsection{\texorpdfstring{11.2. Cas 3.7 :
\texttt{cases/poiseuille.json} (Poiseuille Plan
Stationnaire)}{11.2. Cas 3.7 : cases/poiseuille.json (Poiseuille Plan Stationnaire)}}\label{cas-3.7-casespoiseuille.json-poiseuille-plan-stationnaire}

\begin{Shaded}
\begin{Highlighting}[]
\FunctionTok{\{}
  \DataTypeTok{"case\_name"}\FunctionTok{:} \StringTok{"poiseuille"}\FunctionTok{,}
  \DataTypeTok{"mesh"}\FunctionTok{:} \FunctionTok{\{} \DataTypeTok{"nx"}\FunctionTok{:} \DecValTok{64}\FunctionTok{,} \DataTypeTok{"ny"}\FunctionTok{:} \DecValTok{32}\FunctionTok{,} \DataTypeTok{"lx"}\FunctionTok{:} \FloatTok{2.0}\FunctionTok{,} \DataTypeTok{"ly"}\FunctionTok{:} \FloatTok{1.0} \FunctionTok{\},}
  \DataTypeTok{"physics"}\FunctionTok{:} \FunctionTok{\{}
    \DataTypeTok{"enable\_convection"}\FunctionTok{:} \KeywordTok{false}\FunctionTok{,}
    \DataTypeTok{"enable\_two\_phase"}\FunctionTok{:} \KeywordTok{false}\FunctionTok{,}
    \DataTypeTok{"enable\_surface\_tension"}\FunctionTok{:} \KeywordTok{false}\FunctionTok{,}
    \DataTypeTok{"gx"}\FunctionTok{:} \FloatTok{1.0}\FunctionTok{,} \DataTypeTok{"gy"}\FunctionTok{:} \FloatTok{0.0}\FunctionTok{,}
    \DataTypeTok{"fluid1"}\FunctionTok{:} \FunctionTok{\{} \DataTypeTok{"rho"}\FunctionTok{:} \FloatTok{1.0}\FunctionTok{,} \DataTypeTok{"mu"}\FunctionTok{:} \FloatTok{0.1} \FunctionTok{\}}
  \FunctionTok{\},}
  \DataTypeTok{"time"}\FunctionTok{:} \FunctionTok{\{} \DataTypeTok{"dt"}\FunctionTok{:} \FloatTok{0.05}\FunctionTok{,} \DataTypeTok{"t\_final"}\FunctionTok{:} \FloatTok{15.0}\FunctionTok{,} \DataTypeTok{"adaptive\_dt"}\FunctionTok{:} \KeywordTok{false} \FunctionTok{\},}
  \DataTypeTok{"boundary\_conditions"}\FunctionTok{:} \FunctionTok{\{}
    \DataTypeTok{"left"}\FunctionTok{:}   \FunctionTok{\{} \DataTypeTok{"type"}\FunctionTok{:} \StringTok{"PERIODIC"} \FunctionTok{\},}
    \DataTypeTok{"right"}\FunctionTok{:}  \FunctionTok{\{} \DataTypeTok{"type"}\FunctionTok{:} \StringTok{"PERIODIC"} \FunctionTok{\},}
    \DataTypeTok{"bottom"}\FunctionTok{:} \FunctionTok{\{} \DataTypeTok{"type"}\FunctionTok{:} \StringTok{"DIRICHLET"}\FunctionTok{,} \DataTypeTok{"u"}\FunctionTok{:} \FloatTok{0.0}\FunctionTok{,} \DataTypeTok{"v"}\FunctionTok{:} \FloatTok{0.0} \FunctionTok{\},}
    \DataTypeTok{"top"}\FunctionTok{:}    \FunctionTok{\{} \DataTypeTok{"type"}\FunctionTok{:} \StringTok{"DIRICHLET"}\FunctionTok{,} \DataTypeTok{"u"}\FunctionTok{:} \FloatTok{0.0}\FunctionTok{,} \DataTypeTok{"v"}\FunctionTok{:} \FloatTok{0.0} \FunctionTok{\}}
  \FunctionTok{\},}
  \DataTypeTok{"initial\_condition"}\FunctionTok{:} \FunctionTok{\{} \DataTypeTok{"type"}\FunctionTok{:} \StringTok{"UNIFORM"}\FunctionTok{,} \DataTypeTok{"u"}\FunctionTok{:} \FloatTok{0.0}\FunctionTok{,} \DataTypeTok{"v"}\FunctionTok{:} \FloatTok{0.0} \FunctionTok{\},}
  \DataTypeTok{"solver"}\FunctionTok{:} \FunctionTok{\{}
    \DataTypeTok{"type"}\FunctionTok{:} \StringTok{"AUGMENTED\_LAGRANGIAN"}\FunctionTok{,}
    \DataTypeTok{"linear\_solver"}\FunctionTok{:} \StringTok{"BICGSTAB2"}\FunctionTok{,}
    \DataTypeTok{"precond\_type"}\FunctionTok{:} \StringTok{"JACOBI"}\FunctionTok{,}
    \DataTypeTok{"tol"}\FunctionTok{:} \FloatTok{1.0e{-}7}\FunctionTok{,} \DataTypeTok{"max\_iter"}\FunctionTok{:} \DecValTok{300}\FunctionTok{,} \DataTypeTok{"r\_al"}\FunctionTok{:} \FloatTok{100.0}\FunctionTok{,} \DataTypeTok{"uzawa\_iter"}\FunctionTok{:} \DecValTok{1}
  \FunctionTok{\},}
  \DataTypeTok{"output"}\FunctionTok{:} \FunctionTok{\{} \DataTypeTok{"directory"}\FunctionTok{:} \StringTok{"output\_poiseuille"}\FunctionTok{,} \DataTypeTok{"vtk\_interval"}\FunctionTok{:} \DecValTok{50}\FunctionTok{,} \DataTypeTok{"console\_interval"}\FunctionTok{:} \DecValTok{20} \FunctionTok{\}}
\FunctionTok{\}}
\end{Highlighting}
\end{Shaded}

\subsubsection{\texorpdfstring{11.3. Cas 3.8 :
\texttt{cases/poiseuille\_unsteady.json} (Poiseuille Transitoire de
Szymanski)}{11.3. Cas 3.8 : cases/poiseuille\_unsteady.json (Poiseuille Transitoire de Szymanski)}}\label{cas-3.8-casespoiseuille_unsteady.json-poiseuille-transitoire-de-szymanski}

\begin{Shaded}
\begin{Highlighting}[]
\FunctionTok{\{}
  \DataTypeTok{"case\_name"}\FunctionTok{:} \StringTok{"poiseuille\_unsteady"}\FunctionTok{,}
  \DataTypeTok{"mesh"}\FunctionTok{:} \FunctionTok{\{} \DataTypeTok{"nx"}\FunctionTok{:} \DecValTok{32}\FunctionTok{,} \DataTypeTok{"ny"}\FunctionTok{:} \DecValTok{64}\FunctionTok{,} \DataTypeTok{"lx"}\FunctionTok{:} \FloatTok{0.02}\FunctionTok{,} \DataTypeTok{"ly"}\FunctionTok{:} \FloatTok{0.02} \FunctionTok{\},}
  \DataTypeTok{"physics"}\FunctionTok{:} \FunctionTok{\{}
    \DataTypeTok{"enable\_convection"}\FunctionTok{:} \KeywordTok{false}\FunctionTok{,}
    \DataTypeTok{"enable\_two\_phase"}\FunctionTok{:} \KeywordTok{false}\FunctionTok{,}
    \DataTypeTok{"enable\_surface\_tension"}\FunctionTok{:} \KeywordTok{false}\FunctionTok{,}
    \DataTypeTok{"gx"}\FunctionTok{:} \FloatTok{5.55e{-}5}\FunctionTok{,} \DataTypeTok{"gy"}\FunctionTok{:} \FloatTok{0.0}\FunctionTok{,}
    \DataTypeTok{"fluid1"}\FunctionTok{:} \FunctionTok{\{} \DataTypeTok{"rho"}\FunctionTok{:} \FloatTok{1000.0}\FunctionTok{,} \DataTypeTok{"mu"}\FunctionTok{:} \FloatTok{1.0e{-}3} \FunctionTok{\}}
  \FunctionTok{\},}
  \DataTypeTok{"time"}\FunctionTok{:} \FunctionTok{\{} \DataTypeTok{"dt"}\FunctionTok{:} \FloatTok{0.01}\FunctionTok{,} \DataTypeTok{"t\_final"}\FunctionTok{:} \FloatTok{2.0}\FunctionTok{,} \DataTypeTok{"adaptive\_dt"}\FunctionTok{:} \KeywordTok{false} \FunctionTok{\},}
  \DataTypeTok{"boundary\_conditions"}\FunctionTok{:} \FunctionTok{\{}
    \DataTypeTok{"left"}\FunctionTok{:}   \FunctionTok{\{} \DataTypeTok{"type"}\FunctionTok{:} \StringTok{"PERIODIC"} \FunctionTok{\},}
    \DataTypeTok{"right"}\FunctionTok{:}  \FunctionTok{\{} \DataTypeTok{"type"}\FunctionTok{:} \StringTok{"PERIODIC"} \FunctionTok{\},}
    \DataTypeTok{"bottom"}\FunctionTok{:} \FunctionTok{\{} \DataTypeTok{"type"}\FunctionTok{:} \StringTok{"DIRICHLET"}\FunctionTok{,} \DataTypeTok{"u"}\FunctionTok{:} \FloatTok{0.0}\FunctionTok{,} \DataTypeTok{"v"}\FunctionTok{:} \FloatTok{0.0} \FunctionTok{\},}
    \DataTypeTok{"top"}\FunctionTok{:}    \FunctionTok{\{} \DataTypeTok{"type"}\FunctionTok{:} \StringTok{"DIRICHLET"}\FunctionTok{,} \DataTypeTok{"u"}\FunctionTok{:} \FloatTok{0.0}\FunctionTok{,} \DataTypeTok{"v"}\FunctionTok{:} \FloatTok{0.0} \FunctionTok{\}}
  \FunctionTok{\},}
  \DataTypeTok{"initial\_condition"}\FunctionTok{:} \FunctionTok{\{} \DataTypeTok{"type"}\FunctionTok{:} \StringTok{"UNIFORM"}\FunctionTok{,} \DataTypeTok{"u"}\FunctionTok{:} \FloatTok{0.0}\FunctionTok{,} \DataTypeTok{"v"}\FunctionTok{:} \FloatTok{0.0} \FunctionTok{\},}
  \DataTypeTok{"solver"}\FunctionTok{:} \FunctionTok{\{}
    \DataTypeTok{"type"}\FunctionTok{:} \StringTok{"AUGMENTED\_LAGRANGIAN"}\FunctionTok{,}
    \DataTypeTok{"linear\_solver"}\FunctionTok{:} \StringTok{"BICGSTAB2"}\FunctionTok{,}
    \DataTypeTok{"precond\_type"}\FunctionTok{:} \StringTok{"JACOBI"}\FunctionTok{,}
    \DataTypeTok{"tol"}\FunctionTok{:} \FloatTok{1.0e{-}6}\FunctionTok{,} \DataTypeTok{"max\_iter"}\FunctionTok{:} \DecValTok{300}\FunctionTok{,} \DataTypeTok{"r\_al"}\FunctionTok{:} \FloatTok{1000.0}\FunctionTok{,} \DataTypeTok{"uzawa\_iter"}\FunctionTok{:} \DecValTok{1}
  \FunctionTok{\},}
  \DataTypeTok{"output"}\FunctionTok{:} \FunctionTok{\{} \DataTypeTok{"directory"}\FunctionTok{:} \StringTok{"output\_poiseuille\_unsteady"}\FunctionTok{,} \DataTypeTok{"vtk\_interval"}\FunctionTok{:} \DecValTok{20}\FunctionTok{,} \DataTypeTok{"console\_interval"}\FunctionTok{:} \DecValTok{10} \FunctionTok{\}}
\FunctionTok{\}}
\end{Highlighting}
\end{Shaded}

\subsubsection{\texorpdfstring{11.4. Cas 3.9 :
\texttt{cases/poiseuille\_variable\_viscosity.json} (Viscosité
Exponentielle)}{11.4. Cas 3.9 : cases/poiseuille\_variable\_viscosity.json (Viscosité Exponentielle)}}\label{cas-3.9-casespoiseuille_variable_viscosity.json-viscosituxe9-exponentielle}

\begin{Shaded}
\begin{Highlighting}[]
\FunctionTok{\{}
  \DataTypeTok{"case\_name"}\FunctionTok{:} \StringTok{"poiseuille\_variable\_viscosity"}\FunctionTok{,}
  \DataTypeTok{"mesh"}\FunctionTok{:} \FunctionTok{\{} \DataTypeTok{"nx"}\FunctionTok{:} \DecValTok{32}\FunctionTok{,} \DataTypeTok{"ny"}\FunctionTok{:} \DecValTok{64}\FunctionTok{,} \DataTypeTok{"lx"}\FunctionTok{:} \FloatTok{1.0}\FunctionTok{,} \DataTypeTok{"ly"}\FunctionTok{:} \FloatTok{1.0} \FunctionTok{\},}
  \DataTypeTok{"physics"}\FunctionTok{:} \FunctionTok{\{}
    \DataTypeTok{"enable\_convection"}\FunctionTok{:} \KeywordTok{false}\FunctionTok{,}
    \DataTypeTok{"enable\_two\_phase"}\FunctionTok{:} \KeywordTok{false}\FunctionTok{,}
    \DataTypeTok{"enable\_surface\_tension"}\FunctionTok{:} \KeywordTok{false}\FunctionTok{,}
    \DataTypeTok{"viscosity\_model"}\FunctionTok{:} \StringTok{"EXPONENTIAL"}\FunctionTok{,}
    \DataTypeTok{"viscosity\_exp\_mu0"}\FunctionTok{:} \FloatTok{0.1}\FunctionTok{,}
    \DataTypeTok{"viscosity\_exp\_a"}\FunctionTok{:} \FloatTok{2.0}\FunctionTok{,}
    \DataTypeTok{"viscosity\_exp\_e"}\FunctionTok{:} \FloatTok{0.5}\FunctionTok{,}
    \DataTypeTok{"gx"}\FunctionTok{:} \FloatTok{1.0}\FunctionTok{,} \DataTypeTok{"gy"}\FunctionTok{:} \FloatTok{0.0}\FunctionTok{,}
    \DataTypeTok{"fluid1"}\FunctionTok{:} \FunctionTok{\{} \DataTypeTok{"rho"}\FunctionTok{:} \FloatTok{1.0}\FunctionTok{,} \DataTypeTok{"mu"}\FunctionTok{:} \FloatTok{0.1} \FunctionTok{\}}
  \FunctionTok{\},}
  \DataTypeTok{"time"}\FunctionTok{:} \FunctionTok{\{} \DataTypeTok{"dt"}\FunctionTok{:} \FloatTok{0.05}\FunctionTok{,} \DataTypeTok{"t\_final"}\FunctionTok{:} \FloatTok{10.0}\FunctionTok{,} \DataTypeTok{"adaptive\_dt"}\FunctionTok{:} \KeywordTok{false} \FunctionTok{\},}
  \DataTypeTok{"boundary\_conditions"}\FunctionTok{:} \FunctionTok{\{}
    \DataTypeTok{"left"}\FunctionTok{:}   \FunctionTok{\{} \DataTypeTok{"type"}\FunctionTok{:} \StringTok{"PERIODIC"} \FunctionTok{\},}
    \DataTypeTok{"right"}\FunctionTok{:}  \FunctionTok{\{} \DataTypeTok{"type"}\FunctionTok{:} \StringTok{"PERIODIC"} \FunctionTok{\},}
    \DataTypeTok{"bottom"}\FunctionTok{:} \FunctionTok{\{} \DataTypeTok{"type"}\FunctionTok{:} \StringTok{"DIRICHLET"}\FunctionTok{,} \DataTypeTok{"u"}\FunctionTok{:} \FloatTok{0.0}\FunctionTok{,} \DataTypeTok{"v"}\FunctionTok{:} \FloatTok{0.0} \FunctionTok{\},}
    \DataTypeTok{"top"}\FunctionTok{:}    \FunctionTok{\{} \DataTypeTok{"type"}\FunctionTok{:} \StringTok{"DIRICHLET"}\FunctionTok{,} \DataTypeTok{"u"}\FunctionTok{:} \FloatTok{0.0}\FunctionTok{,} \DataTypeTok{"v"}\FunctionTok{:} \FloatTok{0.0} \FunctionTok{\}}
  \FunctionTok{\},}
  \DataTypeTok{"initial\_condition"}\FunctionTok{:} \FunctionTok{\{} \DataTypeTok{"type"}\FunctionTok{:} \StringTok{"UNIFORM"}\FunctionTok{,} \DataTypeTok{"u"}\FunctionTok{:} \FloatTok{0.0}\FunctionTok{,} \DataTypeTok{"v"}\FunctionTok{:} \FloatTok{0.0} \FunctionTok{\},}
  \DataTypeTok{"solver"}\FunctionTok{:} \FunctionTok{\{}
    \DataTypeTok{"type"}\FunctionTok{:} \StringTok{"AUGMENTED\_LAGRANGIAN"}\FunctionTok{,}
    \DataTypeTok{"linear\_solver"}\FunctionTok{:} \StringTok{"BICGSTAB2"}\FunctionTok{,}
    \DataTypeTok{"precond\_type"}\FunctionTok{:} \StringTok{"JACOBI"}\FunctionTok{,}
    \DataTypeTok{"tol"}\FunctionTok{:} \FloatTok{1.0e{-}6}\FunctionTok{,} \DataTypeTok{"max\_iter"}\FunctionTok{:} \DecValTok{300}\FunctionTok{,} \DataTypeTok{"r\_al"}\FunctionTok{:} \FloatTok{100.0}\FunctionTok{,} \DataTypeTok{"uzawa\_iter"}\FunctionTok{:} \DecValTok{1}
  \FunctionTok{\},}
  \DataTypeTok{"output"}\FunctionTok{:} \FunctionTok{\{} \DataTypeTok{"directory"}\FunctionTok{:} \StringTok{"output\_poiseuille\_variable\_viscosity"}\FunctionTok{,} \DataTypeTok{"vtk\_interval"}\FunctionTok{:} \DecValTok{40}\FunctionTok{,} \DataTypeTok{"console\_interval"}\FunctionTok{:} \DecValTok{20} \FunctionTok{\}}
\FunctionTok{\}}
\end{Highlighting}
\end{Shaded}

\subsubsection{\texorpdfstring{11.5. Cas 3.19 :
\texttt{cases/vortex\_shear\_advection.json} (Vortex avec
Cisaillement)}{11.5. Cas 3.19 : cases/vortex\_shear\_advection.json (Vortex avec Cisaillement)}}\label{cas-3.19-casesvortex_shear_advection.json-vortex-avec-cisaillement}

\begin{Shaded}
\begin{Highlighting}[]
\FunctionTok{\{}
  \DataTypeTok{"case\_name"}\FunctionTok{:} \StringTok{"vortex\_shear"}\FunctionTok{,}
  \DataTypeTok{"mesh"}\FunctionTok{:} \FunctionTok{\{} \DataTypeTok{"nx"}\FunctionTok{:} \DecValTok{64}\FunctionTok{,} \DataTypeTok{"ny"}\FunctionTok{:} \DecValTok{64}\FunctionTok{,} \DataTypeTok{"lx"}\FunctionTok{:} \FloatTok{1.0}\FunctionTok{,} \DataTypeTok{"ly"}\FunctionTok{:} \FloatTok{1.0} \FunctionTok{\},}
  \DataTypeTok{"physics"}\FunctionTok{:} \FunctionTok{\{}
    \DataTypeTok{"enable\_convection"}\FunctionTok{:} \KeywordTok{false}\FunctionTok{,}
    \DataTypeTok{"enable\_two\_phase"}\FunctionTok{:} \KeywordTok{true}\FunctionTok{,}
    \DataTypeTok{"enable\_surface\_tension"}\FunctionTok{:} \KeywordTok{false}\FunctionTok{,}
    \DataTypeTok{"convective\_scheme"}\FunctionTok{:} \StringTok{"TVD\_SUPERBEE"}\FunctionTok{,}
    \DataTypeTok{"gx"}\FunctionTok{:} \FloatTok{0.0}\FunctionTok{,} \DataTypeTok{"gy"}\FunctionTok{:} \FloatTok{0.0}\FunctionTok{,}
    \DataTypeTok{"fluid1"}\FunctionTok{:} \FunctionTok{\{} \DataTypeTok{"rho"}\FunctionTok{:} \FloatTok{1.0}\FunctionTok{,} \DataTypeTok{"mu"}\FunctionTok{:} \FloatTok{1.0} \FunctionTok{\},}
    \DataTypeTok{"fluid2"}\FunctionTok{:} \FunctionTok{\{} \DataTypeTok{"rho"}\FunctionTok{:} \FloatTok{1.0}\FunctionTok{,} \DataTypeTok{"mu"}\FunctionTok{:} \FloatTok{1.0} \FunctionTok{\}}
  \FunctionTok{\},}
  \DataTypeTok{"time"}\FunctionTok{:} \FunctionTok{\{} \DataTypeTok{"dt"}\FunctionTok{:} \FloatTok{0.005}\FunctionTok{,} \DataTypeTok{"t\_final"}\FunctionTok{:} \FloatTok{4.0}\FunctionTok{,} \DataTypeTok{"adaptive\_dt"}\FunctionTok{:} \KeywordTok{false} \FunctionTok{\},}
  \DataTypeTok{"boundary\_conditions"}\FunctionTok{:} \FunctionTok{\{}
    \DataTypeTok{"left"}\FunctionTok{:}   \FunctionTok{\{} \DataTypeTok{"type"}\FunctionTok{:} \StringTok{"FREE\_SLIP"} \FunctionTok{\},}
    \DataTypeTok{"right"}\FunctionTok{:}  \FunctionTok{\{} \DataTypeTok{"type"}\FunctionTok{:} \StringTok{"FREE\_SLIP"} \FunctionTok{\},}
    \DataTypeTok{"bottom"}\FunctionTok{:} \FunctionTok{\{} \DataTypeTok{"type"}\FunctionTok{:} \StringTok{"FREE\_SLIP"} \FunctionTok{\},}
    \DataTypeTok{"top"}\FunctionTok{:}    \FunctionTok{\{} \DataTypeTok{"type"}\FunctionTok{:} \StringTok{"FREE\_SLIP"} \FunctionTok{\}}
  \FunctionTok{\},}
  \DataTypeTok{"initial\_condition"}\FunctionTok{:} \FunctionTok{\{}
    \DataTypeTok{"type"}\FunctionTok{:} \StringTok{"VORTEX\_SHEAR"}\FunctionTok{,}
    \DataTypeTok{"center\_x"}\FunctionTok{:} \FloatTok{0.5}\FunctionTok{,} \DataTypeTok{"center\_y"}\FunctionTok{:} \FloatTok{0.75}\FunctionTok{,} \DataTypeTok{"radius"}\FunctionTok{:} \FloatTok{0.15}
  \FunctionTok{\},}
  \DataTypeTok{"solver"}\FunctionTok{:} \FunctionTok{\{} \DataTypeTok{"type"}\FunctionTok{:} \StringTok{"AUGMENTED\_LAGRANGIAN"} \FunctionTok{\},}
  \DataTypeTok{"output"}\FunctionTok{:} \FunctionTok{\{} \DataTypeTok{"directory"}\FunctionTok{:} \StringTok{"output\_vortex\_shear"}\FunctionTok{,} \DataTypeTok{"vtk\_interval"}\FunctionTok{:} \DecValTok{50}\FunctionTok{,} \DataTypeTok{"console\_interval"}\FunctionTok{:} \DecValTok{50} \FunctionTok{\}}
\FunctionTok{\}}
\end{Highlighting}
\end{Shaded}

\subsubsection{\texorpdfstring{11.6. Cas 4.1 :
\texttt{cases/backward\_facing\_step.json} (Marche Descendante
\(Re=100\))}{11.6. Cas 4.1 : cases/backward\_facing\_step.json (Marche Descendante Re=100)}}\label{cas-4.1-casesbackward_facing_step.json-marche-descendante-re100}

\begin{Shaded}
\begin{Highlighting}[]
\FunctionTok{\{}
  \DataTypeTok{"case\_name"}\FunctionTok{:} \StringTok{"backward\_facing\_step"}\FunctionTok{,}
  \DataTypeTok{"mesh"}\FunctionTok{:} \FunctionTok{\{} \DataTypeTok{"nx"}\FunctionTok{:} \DecValTok{128}\FunctionTok{,} \DataTypeTok{"ny"}\FunctionTok{:} \DecValTok{32}\FunctionTok{,} \DataTypeTok{"lx"}\FunctionTok{:} \FloatTok{40.0}\FunctionTok{,} \DataTypeTok{"ly"}\FunctionTok{:} \FloatTok{2.0612} \FunctionTok{\},}
  \DataTypeTok{"physics"}\FunctionTok{:} \FunctionTok{\{}
    \DataTypeTok{"enable\_convection"}\FunctionTok{:} \KeywordTok{true}\FunctionTok{,}
    \DataTypeTok{"enable\_two\_phase"}\FunctionTok{:} \KeywordTok{false}\FunctionTok{,}
    \DataTypeTok{"enable\_surface\_tension"}\FunctionTok{:} \KeywordTok{false}\FunctionTok{,}
    \DataTypeTok{"gx"}\FunctionTok{:} \FloatTok{0.0}\FunctionTok{,} \DataTypeTok{"gy"}\FunctionTok{:} \FloatTok{0.0}\FunctionTok{,}
    \DataTypeTok{"fluid1"}\FunctionTok{:} \FunctionTok{\{} \DataTypeTok{"rho"}\FunctionTok{:} \FloatTok{1.0}\FunctionTok{,} \DataTypeTok{"mu"}\FunctionTok{:} \FloatTok{0.01} \FunctionTok{\}}
  \FunctionTok{\},}
  \DataTypeTok{"immersed\_boundary"}\FunctionTok{:} \FunctionTok{\{}
    \DataTypeTok{"enabled"}\FunctionTok{:} \KeywordTok{true}\FunctionTok{,}
    \DataTypeTok{"type"}\FunctionTok{:} \StringTok{"BOX"}\FunctionTok{,}
    \DataTypeTok{"box\_x0"}\FunctionTok{:} \FloatTok{0.0}\FunctionTok{,} \DataTypeTok{"box\_y0"}\FunctionTok{:} \FloatTok{0.0}\FunctionTok{,} \DataTypeTok{"box\_x1"}\FunctionTok{:} \FloatTok{4.0}\FunctionTok{,} \DataTypeTok{"box\_y1"}\FunctionTok{:} \FloatTok{1.0}\FunctionTok{,}
    \DataTypeTok{"permeability\_solid"}\FunctionTok{:} \FloatTok{1.0e{-}40}
  \FunctionTok{\},}
  \DataTypeTok{"time"}\FunctionTok{:} \FunctionTok{\{} \DataTypeTok{"dt"}\FunctionTok{:} \FloatTok{0.05}\FunctionTok{,} \DataTypeTok{"t\_final"}\FunctionTok{:} \FloatTok{60.0}\FunctionTok{,} \DataTypeTok{"adaptive\_dt"}\FunctionTok{:} \KeywordTok{false} \FunctionTok{\},}
  \DataTypeTok{"boundary\_conditions"}\FunctionTok{:} \FunctionTok{\{}
    \DataTypeTok{"left"}\FunctionTok{:}   \FunctionTok{\{} \DataTypeTok{"type"}\FunctionTok{:} \StringTok{"DIRICHLET"}\FunctionTok{,} \DataTypeTok{"u"}\FunctionTok{:} \FloatTok{1.0}\FunctionTok{,} \DataTypeTok{"v"}\FunctionTok{:} \FloatTok{0.0} \FunctionTok{\},}
    \DataTypeTok{"right"}\FunctionTok{:}  \FunctionTok{\{} \DataTypeTok{"type"}\FunctionTok{:} \StringTok{"NEUMANN"} \FunctionTok{\},}
    \DataTypeTok{"bottom"}\FunctionTok{:} \FunctionTok{\{} \DataTypeTok{"type"}\FunctionTok{:} \StringTok{"DIRICHLET"}\FunctionTok{,} \DataTypeTok{"u"}\FunctionTok{:} \FloatTok{0.0}\FunctionTok{,} \DataTypeTok{"v"}\FunctionTok{:} \FloatTok{0.0} \FunctionTok{\},}
    \DataTypeTok{"top"}\FunctionTok{:}    \FunctionTok{\{} \DataTypeTok{"type"}\FunctionTok{:} \StringTok{"DIRICHLET"}\FunctionTok{,} \DataTypeTok{"u"}\FunctionTok{:} \FloatTok{0.0}\FunctionTok{,} \DataTypeTok{"v"}\FunctionTok{:} \FloatTok{0.0} \FunctionTok{\}}
  \FunctionTok{\},}
  \DataTypeTok{"initial\_condition"}\FunctionTok{:} \FunctionTok{\{} \DataTypeTok{"type"}\FunctionTok{:} \StringTok{"UNIFORM"}\FunctionTok{,} \DataTypeTok{"u"}\FunctionTok{:} \FloatTok{0.0}\FunctionTok{,} \DataTypeTok{"v"}\FunctionTok{:} \FloatTok{0.0} \FunctionTok{\},}
  \DataTypeTok{"solver"}\FunctionTok{:} \FunctionTok{\{}
    \DataTypeTok{"type"}\FunctionTok{:} \StringTok{"AUGMENTED\_LAGRANGIAN"}\FunctionTok{,}
    \DataTypeTok{"linear\_solver"}\FunctionTok{:} \StringTok{"BICGSTAB2"}\FunctionTok{,}
    \DataTypeTok{"precond\_type"}\FunctionTok{:} \StringTok{"JACOBI"}\FunctionTok{,}
    \DataTypeTok{"tol"}\FunctionTok{:} \FloatTok{1.0e{-}5}\FunctionTok{,} \DataTypeTok{"max\_iter"}\FunctionTok{:} \DecValTok{300}\FunctionTok{,} \DataTypeTok{"r\_al"}\FunctionTok{:} \FloatTok{100.0}\FunctionTok{,} \DataTypeTok{"uzawa\_iter"}\FunctionTok{:} \DecValTok{1}
  \FunctionTok{\},}
  \DataTypeTok{"output"}\FunctionTok{:} \FunctionTok{\{} \DataTypeTok{"directory"}\FunctionTok{:} \StringTok{"output\_step"}\FunctionTok{,} \DataTypeTok{"vtk\_interval"}\FunctionTok{:} \DecValTok{100}\FunctionTok{,} \DataTypeTok{"console\_interval"}\FunctionTok{:} \DecValTok{50} \FunctionTok{\}}
\FunctionTok{\}}
\end{Highlighting}
\end{Shaded}

\subsubsection{\texorpdfstring{11.7. Cas 4.4 :
\texttt{cases/obstacle\_cylinder\_re20.json} (Cylindre Immergé
\(Re=20\))}{11.7. Cas 4.4 : cases/obstacle\_cylinder\_re20.json (Cylindre Immergé Re=20)}}\label{cas-4.4-casesobstacle_cylinder_re20.json-cylindre-immerguxe9-re20}

\begin{Shaded}
\begin{Highlighting}[]
\FunctionTok{\{}
  \DataTypeTok{"case\_name"}\FunctionTok{:} \StringTok{"cylinder\_re20"}\FunctionTok{,}
  \DataTypeTok{"mesh"}\FunctionTok{:} \FunctionTok{\{} \DataTypeTok{"nx"}\FunctionTok{:} \DecValTok{64}\FunctionTok{,} \DataTypeTok{"ny"}\FunctionTok{:} \DecValTok{64}\FunctionTok{,} \DataTypeTok{"lx"}\FunctionTok{:} \FloatTok{1.0}\FunctionTok{,} \DataTypeTok{"ly"}\FunctionTok{:} \FloatTok{1.0} \FunctionTok{\},}
  \DataTypeTok{"physics"}\FunctionTok{:} \FunctionTok{\{}
    \DataTypeTok{"enable\_convection"}\FunctionTok{:} \KeywordTok{true}\FunctionTok{,}
    \DataTypeTok{"enable\_two\_phase"}\FunctionTok{:} \KeywordTok{false}\FunctionTok{,}
    \DataTypeTok{"enable\_surface\_tension"}\FunctionTok{:} \KeywordTok{false}\FunctionTok{,}
    \DataTypeTok{"gx"}\FunctionTok{:} \FloatTok{0.0}\FunctionTok{,} \DataTypeTok{"gy"}\FunctionTok{:} \FloatTok{0.0}\FunctionTok{,}
    \DataTypeTok{"fluid1"}\FunctionTok{:} \FunctionTok{\{} \DataTypeTok{"rho"}\FunctionTok{:} \FloatTok{1.0}\FunctionTok{,} \DataTypeTok{"mu"}\FunctionTok{:} \FloatTok{0.01} \FunctionTok{\}}
  \FunctionTok{\},}
  \DataTypeTok{"immersed\_boundary"}\FunctionTok{:} \FunctionTok{\{}
    \DataTypeTok{"enabled"}\FunctionTok{:} \KeywordTok{true}\FunctionTok{,}
    \DataTypeTok{"type"}\FunctionTok{:} \StringTok{"CYLINDER"}\FunctionTok{,}
    \DataTypeTok{"cylinder\_xc"}\FunctionTok{:} \FloatTok{0.5}\FunctionTok{,} \DataTypeTok{"cylinder\_yc"}\FunctionTok{:} \FloatTok{0.5}\FunctionTok{,} \DataTypeTok{"cylinder\_r"}\FunctionTok{:} \FloatTok{0.1}\FunctionTok{,}
    \DataTypeTok{"permeability\_solid"}\FunctionTok{:} \FloatTok{1.0e{-}40}
  \FunctionTok{\},}
  \DataTypeTok{"time"}\FunctionTok{:} \FunctionTok{\{} \DataTypeTok{"dt"}\FunctionTok{:} \FloatTok{0.05}\FunctionTok{,} \DataTypeTok{"t\_final"}\FunctionTok{:} \FloatTok{15.0}\FunctionTok{,} \DataTypeTok{"adaptive\_dt"}\FunctionTok{:} \KeywordTok{false} \FunctionTok{\},}
  \DataTypeTok{"boundary\_conditions"}\FunctionTok{:} \FunctionTok{\{}
    \DataTypeTok{"left"}\FunctionTok{:}   \FunctionTok{\{} \DataTypeTok{"type"}\FunctionTok{:} \StringTok{"DIRICHLET"}\FunctionTok{,} \DataTypeTok{"u"}\FunctionTok{:} \FloatTok{1.0}\FunctionTok{,} \DataTypeTok{"v"}\FunctionTok{:} \FloatTok{0.0} \FunctionTok{\},}
    \DataTypeTok{"right"}\FunctionTok{:}  \FunctionTok{\{} \DataTypeTok{"type"}\FunctionTok{:} \StringTok{"NEUMANN"} \FunctionTok{\},}
    \DataTypeTok{"bottom"}\FunctionTok{:} \FunctionTok{\{} \DataTypeTok{"type"}\FunctionTok{:} \StringTok{"FREE\_SLIP"} \FunctionTok{\},}
    \DataTypeTok{"top"}\FunctionTok{:}    \FunctionTok{\{} \DataTypeTok{"type"}\FunctionTok{:} \StringTok{"FREE\_SLIP"} \FunctionTok{\}}
  \FunctionTok{\},}
  \DataTypeTok{"initial\_condition"}\FunctionTok{:} \FunctionTok{\{} \DataTypeTok{"type"}\FunctionTok{:} \StringTok{"UNIFORM"}\FunctionTok{,} \DataTypeTok{"u"}\FunctionTok{:} \FloatTok{1.0}\FunctionTok{,} \DataTypeTok{"v"}\FunctionTok{:} \FloatTok{0.0} \FunctionTok{\},}
  \DataTypeTok{"solver"}\FunctionTok{:} \FunctionTok{\{}
    \DataTypeTok{"type"}\FunctionTok{:} \StringTok{"AUGMENTED\_LAGRANGIAN"}\FunctionTok{,}
    \DataTypeTok{"linear\_solver"}\FunctionTok{:} \StringTok{"BICGSTAB2"}\FunctionTok{,}
    \DataTypeTok{"precond\_type"}\FunctionTok{:} \StringTok{"JACOBI"}\FunctionTok{,}
    \DataTypeTok{"tol"}\FunctionTok{:} \FloatTok{1.0e{-}5}\FunctionTok{,} \DataTypeTok{"max\_iter"}\FunctionTok{:} \DecValTok{300}\FunctionTok{,} \DataTypeTok{"r\_al"}\FunctionTok{:} \FloatTok{100.0}\FunctionTok{,} \DataTypeTok{"uzawa\_iter"}\FunctionTok{:} \DecValTok{1}
  \FunctionTok{\},}
  \DataTypeTok{"output"}\FunctionTok{:} \FunctionTok{\{} \DataTypeTok{"directory"}\FunctionTok{:} \StringTok{"output\_cylinder\_re20"}\FunctionTok{,} \DataTypeTok{"vtk\_interval"}\FunctionTok{:} \DecValTok{50}\FunctionTok{,} \DataTypeTok{"console\_interval"}\FunctionTok{:} \DecValTok{20} \FunctionTok{\}}
\FunctionTok{\}}
\end{Highlighting}
\end{Shaded}

\subsubsection{\texorpdfstring{11.8. Cas 4.5 :
\texttt{cases/cavity\_re1000.json} (Cavité Carrée Entraînée
\(Re=1000\))}{11.8. Cas 4.5 : cases/cavity\_re1000.json (Cavité Carrée Entraînée Re=1000)}}\label{cas-4.5-casescavity_re1000.json-cavituxe9-carruxe9e-entrauxeenuxe9e-re1000}

\begin{Shaded}
\begin{Highlighting}[]
\FunctionTok{\{}
  \DataTypeTok{"case\_name"}\FunctionTok{:} \StringTok{"cavity\_re1000"}\FunctionTok{,}
  \DataTypeTok{"mesh"}\FunctionTok{:} \FunctionTok{\{} \DataTypeTok{"nx"}\FunctionTok{:} \DecValTok{128}\FunctionTok{,} \DataTypeTok{"ny"}\FunctionTok{:} \DecValTok{128}\FunctionTok{,} \DataTypeTok{"lx"}\FunctionTok{:} \FloatTok{1.0}\FunctionTok{,} \DataTypeTok{"ly"}\FunctionTok{:} \FloatTok{1.0} \FunctionTok{\},}
  \DataTypeTok{"physics"}\FunctionTok{:} \FunctionTok{\{}
    \DataTypeTok{"enable\_convection"}\FunctionTok{:} \KeywordTok{true}\FunctionTok{,}
    \DataTypeTok{"enable\_two\_phase"}\FunctionTok{:} \KeywordTok{false}\FunctionTok{,}
    \DataTypeTok{"enable\_surface\_tension"}\FunctionTok{:} \KeywordTok{false}\FunctionTok{,}
    \DataTypeTok{"gx"}\FunctionTok{:} \FloatTok{0.0}\FunctionTok{,} \DataTypeTok{"gy"}\FunctionTok{:} \FloatTok{0.0}\FunctionTok{,}
    \DataTypeTok{"fluid1"}\FunctionTok{:} \FunctionTok{\{} \DataTypeTok{"rho"}\FunctionTok{:} \FloatTok{1.0}\FunctionTok{,} \DataTypeTok{"mu"}\FunctionTok{:} \FloatTok{0.001} \FunctionTok{\}}
  \FunctionTok{\},}
  \DataTypeTok{"time"}\FunctionTok{:} \FunctionTok{\{} \DataTypeTok{"dt"}\FunctionTok{:} \FloatTok{0.1}\FunctionTok{,} \DataTypeTok{"t\_final"}\FunctionTok{:} \FloatTok{80.0}\FunctionTok{,} \DataTypeTok{"adaptive\_dt"}\FunctionTok{:} \KeywordTok{false} \FunctionTok{\},}
  \DataTypeTok{"boundary\_conditions"}\FunctionTok{:} \FunctionTok{\{}
    \DataTypeTok{"left"}\FunctionTok{:}   \FunctionTok{\{} \DataTypeTok{"type"}\FunctionTok{:} \StringTok{"DIRICHLET"}\FunctionTok{,} \DataTypeTok{"u"}\FunctionTok{:} \FloatTok{0.0}\FunctionTok{,} \DataTypeTok{"v"}\FunctionTok{:} \FloatTok{0.0} \FunctionTok{\},}
    \DataTypeTok{"right"}\FunctionTok{:}  \FunctionTok{\{} \DataTypeTok{"type"}\FunctionTok{:} \StringTok{"DIRICHLET"}\FunctionTok{,} \DataTypeTok{"u"}\FunctionTok{:} \FloatTok{0.0}\FunctionTok{,} \DataTypeTok{"v"}\FunctionTok{:} \FloatTok{0.0} \FunctionTok{\},}
    \DataTypeTok{"bottom"}\FunctionTok{:} \FunctionTok{\{} \DataTypeTok{"type"}\FunctionTok{:} \StringTok{"DIRICHLET"}\FunctionTok{,} \DataTypeTok{"u"}\FunctionTok{:} \FloatTok{0.0}\FunctionTok{,} \DataTypeTok{"v"}\FunctionTok{:} \FloatTok{0.0} \FunctionTok{\},}
    \DataTypeTok{"top"}\FunctionTok{:}    \FunctionTok{\{} \DataTypeTok{"type"}\FunctionTok{:} \StringTok{"DIRICHLET"}\FunctionTok{,} \DataTypeTok{"u"}\FunctionTok{:} \FloatTok{1.0}\FunctionTok{,} \DataTypeTok{"v"}\FunctionTok{:} \FloatTok{0.0} \FunctionTok{\}}
  \FunctionTok{\},}
  \DataTypeTok{"initial\_condition"}\FunctionTok{:} \FunctionTok{\{} \DataTypeTok{"type"}\FunctionTok{:} \StringTok{"UNIFORM"}\FunctionTok{,} \DataTypeTok{"u"}\FunctionTok{:} \FloatTok{0.0}\FunctionTok{,} \DataTypeTok{"v"}\FunctionTok{:} \FloatTok{0.0} \FunctionTok{\},}
  \DataTypeTok{"solver"}\FunctionTok{:} \FunctionTok{\{}
    \DataTypeTok{"type"}\FunctionTok{:} \StringTok{"AUGMENTED\_LAGRANGIAN"}\FunctionTok{,}
    \DataTypeTok{"linear\_solver"}\FunctionTok{:} \StringTok{"BICGSTAB2"}\FunctionTok{,}
    \DataTypeTok{"precond\_type"}\FunctionTok{:} \StringTok{"JACOBI"}\FunctionTok{,}
    \DataTypeTok{"tol"}\FunctionTok{:} \FloatTok{1.0e{-}5}\FunctionTok{,} \DataTypeTok{"max\_iter"}\FunctionTok{:} \DecValTok{300}\FunctionTok{,} \DataTypeTok{"r\_al"}\FunctionTok{:} \FloatTok{100.0}\FunctionTok{,} \DataTypeTok{"uzawa\_iter"}\FunctionTok{:} \DecValTok{1}
  \FunctionTok{\},}
  \DataTypeTok{"output"}\FunctionTok{:} \FunctionTok{\{} \DataTypeTok{"directory"}\FunctionTok{:} \StringTok{"output\_cavity\_re1000"}\FunctionTok{,} \DataTypeTok{"vtk\_interval"}\FunctionTok{:} \DecValTok{200}\FunctionTok{,} \DataTypeTok{"console\_interval"}\FunctionTok{:} \DecValTok{50} \FunctionTok{\}}
\FunctionTok{\}}
\end{Highlighting}
\end{Shaded}

\subsubsection{\texorpdfstring{11.9. Cas 4.7 :
\texttt{cases/cavity\_l\_shaped.json} (Cavité en Forme de L
\(Re=1000\))}{11.9. Cas 4.7 : cases/cavity\_l\_shaped.json (Cavité en Forme de L Re=1000)}}\label{cas-4.7-casescavity_l_shaped.json-cavituxe9-en-forme-de-l-re1000}

\begin{Shaded}
\begin{Highlighting}[]
\FunctionTok{\{}
  \DataTypeTok{"case\_name"}\FunctionTok{:} \StringTok{"cavity\_l\_shaped"}\FunctionTok{,}
  \DataTypeTok{"mesh"}\FunctionTok{:} \FunctionTok{\{} \DataTypeTok{"nx"}\FunctionTok{:} \DecValTok{64}\FunctionTok{,} \DataTypeTok{"ny"}\FunctionTok{:} \DecValTok{64}\FunctionTok{,} \DataTypeTok{"lx"}\FunctionTok{:} \FloatTok{1.0}\FunctionTok{,} \DataTypeTok{"ly"}\FunctionTok{:} \FloatTok{1.0} \FunctionTok{\},}
  \DataTypeTok{"physics"}\FunctionTok{:} \FunctionTok{\{}
    \DataTypeTok{"enable\_convection"}\FunctionTok{:} \KeywordTok{true}\FunctionTok{,}
    \DataTypeTok{"enable\_two\_phase"}\FunctionTok{:} \KeywordTok{false}\FunctionTok{,}
    \DataTypeTok{"enable\_surface\_tension"}\FunctionTok{:} \KeywordTok{false}\FunctionTok{,}
    \DataTypeTok{"gx"}\FunctionTok{:} \FloatTok{0.0}\FunctionTok{,} \DataTypeTok{"gy"}\FunctionTok{:} \FloatTok{0.0}\FunctionTok{,}
    \DataTypeTok{"fluid1"}\FunctionTok{:} \FunctionTok{\{} \DataTypeTok{"rho"}\FunctionTok{:} \FloatTok{1.0}\FunctionTok{,} \DataTypeTok{"mu"}\FunctionTok{:} \FloatTok{0.001} \FunctionTok{\}}
  \FunctionTok{\},}
  \DataTypeTok{"immersed\_boundary"}\FunctionTok{:} \FunctionTok{\{}
    \DataTypeTok{"enabled"}\FunctionTok{:} \KeywordTok{true}\FunctionTok{,}
    \DataTypeTok{"type"}\FunctionTok{:} \StringTok{"BOX"}\FunctionTok{,}
    \DataTypeTok{"box\_x0"}\FunctionTok{:} \FloatTok{0.5}\FunctionTok{,} \DataTypeTok{"box\_y0"}\FunctionTok{:} \FloatTok{0.5}\FunctionTok{,} \DataTypeTok{"box\_x1"}\FunctionTok{:} \FloatTok{1.0}\FunctionTok{,} \DataTypeTok{"box\_y1"}\FunctionTok{:} \FloatTok{1.0}\FunctionTok{,}
    \DataTypeTok{"permeability\_solid"}\FunctionTok{:} \FloatTok{1.0e{-}40}
  \FunctionTok{\},}
  \DataTypeTok{"time"}\FunctionTok{:} \FunctionTok{\{} \DataTypeTok{"dt"}\FunctionTok{:} \FloatTok{0.05}\FunctionTok{,} \DataTypeTok{"t\_final"}\FunctionTok{:} \FloatTok{15.0}\FunctionTok{,} \DataTypeTok{"adaptive\_dt"}\FunctionTok{:} \KeywordTok{false} \FunctionTok{\},}
  \DataTypeTok{"boundary\_conditions"}\FunctionTok{:} \FunctionTok{\{}
    \DataTypeTok{"left"}\FunctionTok{:}   \FunctionTok{\{} \DataTypeTok{"type"}\FunctionTok{:} \StringTok{"DIRICHLET"}\FunctionTok{,} \DataTypeTok{"u"}\FunctionTok{:} \FloatTok{0.0}\FunctionTok{,} \DataTypeTok{"v"}\FunctionTok{:} \FloatTok{0.0} \FunctionTok{\},}
    \DataTypeTok{"right"}\FunctionTok{:}  \FunctionTok{\{} \DataTypeTok{"type"}\FunctionTok{:} \StringTok{"DIRICHLET"}\FunctionTok{,} \DataTypeTok{"u"}\FunctionTok{:} \FloatTok{0.0}\FunctionTok{,} \DataTypeTok{"v"}\FunctionTok{:} \FloatTok{0.0} \FunctionTok{\},}
    \DataTypeTok{"bottom"}\FunctionTok{:} \FunctionTok{\{} \DataTypeTok{"type"}\FunctionTok{:} \StringTok{"DIRICHLET"}\FunctionTok{,} \DataTypeTok{"u"}\FunctionTok{:} \FloatTok{0.0}\FunctionTok{,} \DataTypeTok{"v"}\FunctionTok{:} \FloatTok{0.0} \FunctionTok{\},}
    \DataTypeTok{"top"}\FunctionTok{:}    \FunctionTok{\{} \DataTypeTok{"type"}\FunctionTok{:} \StringTok{"DIRICHLET"}\FunctionTok{,} \DataTypeTok{"u"}\FunctionTok{:} \FloatTok{1.0}\FunctionTok{,} \DataTypeTok{"v"}\FunctionTok{:} \FloatTok{0.0} \FunctionTok{\}}
  \FunctionTok{\},}
  \DataTypeTok{"initial\_condition"}\FunctionTok{:} \FunctionTok{\{} \DataTypeTok{"type"}\FunctionTok{:} \StringTok{"UNIFORM"}\FunctionTok{,} \DataTypeTok{"u"}\FunctionTok{:} \FloatTok{0.0}\FunctionTok{,} \DataTypeTok{"v"}\FunctionTok{:} \FloatTok{0.0} \FunctionTok{\},}
  \DataTypeTok{"solver"}\FunctionTok{:} \FunctionTok{\{}
    \DataTypeTok{"type"}\FunctionTok{:} \StringTok{"AUGMENTED\_LAGRANGIAN"}\FunctionTok{,}
    \DataTypeTok{"linear\_solver"}\FunctionTok{:} \StringTok{"BICGSTAB2"}\FunctionTok{,}
    \DataTypeTok{"precond\_type"}\FunctionTok{:} \StringTok{"JACOBI"}\FunctionTok{,}
    \DataTypeTok{"tol"}\FunctionTok{:} \FloatTok{1.0e{-}5}\FunctionTok{,} \DataTypeTok{"max\_iter"}\FunctionTok{:} \DecValTok{300}\FunctionTok{,} \DataTypeTok{"r\_al"}\FunctionTok{:} \FloatTok{100.0}\FunctionTok{,} \DataTypeTok{"uzawa\_iter"}\FunctionTok{:} \DecValTok{1}
  \FunctionTok{\},}
  \DataTypeTok{"output"}\FunctionTok{:} \FunctionTok{\{} \DataTypeTok{"directory"}\FunctionTok{:} \StringTok{"output\_cavity\_l"}\FunctionTok{,} \DataTypeTok{"vtk\_interval"}\FunctionTok{:} \DecValTok{50}\FunctionTok{,} \DataTypeTok{"console\_interval"}\FunctionTok{:} \DecValTok{20} \FunctionTok{\}}
\FunctionTok{\}}
\end{Highlighting}
\end{Shaded}

\subsubsection{\texorpdfstring{11.10. Cas 6.1 :
\texttt{cases/natural\_convection\_ra1e4.json} (Convection Naturelle
\(Ra=10^4\))}{11.10. Cas 6.1 : cases/natural\_convection\_ra1e4.json (Convection Naturelle Ra=10\^{}4)}}\label{cas-6.1-casesnatural_convection_ra1e4.json-convection-naturelle-ra104}

\begin{Shaded}
\begin{Highlighting}[]
\FunctionTok{\{}
  \DataTypeTok{"case\_name"}\FunctionTok{:} \StringTok{"natural\_convection\_ra1e4"}\FunctionTok{,}
  \DataTypeTok{"mesh"}\FunctionTok{:} \FunctionTok{\{} \DataTypeTok{"nx"}\FunctionTok{:} \DecValTok{64}\FunctionTok{,} \DataTypeTok{"ny"}\FunctionTok{:} \DecValTok{64}\FunctionTok{,} \DataTypeTok{"lx"}\FunctionTok{:} \FloatTok{1.0}\FunctionTok{,} \DataTypeTok{"ly"}\FunctionTok{:} \FloatTok{1.0} \FunctionTok{\},}
  \DataTypeTok{"physics"}\FunctionTok{:} \FunctionTok{\{}
    \DataTypeTok{"enable\_convection"}\FunctionTok{:} \KeywordTok{true}\FunctionTok{,}
    \DataTypeTok{"enable\_two\_phase"}\FunctionTok{:} \KeywordTok{false}\FunctionTok{,}
    \DataTypeTok{"enable\_surface\_tension"}\FunctionTok{:} \KeywordTok{false}\FunctionTok{,}
    \DataTypeTok{"gx"}\FunctionTok{:} \FloatTok{0.0}\FunctionTok{,} \DataTypeTok{"gy"}\FunctionTok{:} \FloatTok{{-}9.81}\FunctionTok{,}
    \DataTypeTok{"fluid1"}\FunctionTok{:} \FunctionTok{\{} \DataTypeTok{"rho"}\FunctionTok{:} \FloatTok{1.0}\FunctionTok{,} \DataTypeTok{"mu"}\FunctionTok{:} \FloatTok{0.01} \FunctionTok{\}}
  \FunctionTok{\},}
  \DataTypeTok{"thermal"}\FunctionTok{:} \FunctionTok{\{}
    \DataTypeTok{"enable\_energy"}\FunctionTok{:} \KeywordTok{true}\FunctionTok{,}
    \DataTypeTok{"prandtl"}\FunctionTok{:} \FloatTok{0.71}\FunctionTok{,}
    \DataTypeTok{"rayleigh"}\FunctionTok{:} \FloatTok{1.0e4}\FunctionTok{,}
    \DataTypeTok{"beta"}\FunctionTok{:} \FloatTok{0.00341}\FunctionTok{,}
    \DataTypeTok{"t\_ref"}\FunctionTok{:} \FloatTok{0.0}\FunctionTok{,}
    \DataTypeTok{"t\_hot"}\FunctionTok{:} \FloatTok{1.0}\FunctionTok{,}
    \DataTypeTok{"t\_cold"}\FunctionTok{:} \FloatTok{0.0}\FunctionTok{,}
    \DataTypeTok{"k\_thermal"}\FunctionTok{:} \FloatTok{0.01408}\FunctionTok{,}
    \DataTypeTok{"cp\_thermal"}\FunctionTok{:} \FloatTok{1.0}
  \FunctionTok{\},}
  \DataTypeTok{"time"}\FunctionTok{:} \FunctionTok{\{} \DataTypeTok{"dt"}\FunctionTok{:} \FloatTok{0.05}\FunctionTok{,} \DataTypeTok{"t\_final"}\FunctionTok{:} \FloatTok{20.0}\FunctionTok{,} \DataTypeTok{"adaptive\_dt"}\FunctionTok{:} \KeywordTok{false} \FunctionTok{\},}
  \DataTypeTok{"boundary\_conditions"}\FunctionTok{:} \FunctionTok{\{}
    \DataTypeTok{"left"}\FunctionTok{:}   \FunctionTok{\{} \DataTypeTok{"type"}\FunctionTok{:} \StringTok{"DIRICHLET"}\FunctionTok{,} \DataTypeTok{"u"}\FunctionTok{:} \FloatTok{0.0}\FunctionTok{,} \DataTypeTok{"v"}\FunctionTok{:} \FloatTok{0.0}\FunctionTok{,} \DataTypeTok{"thermal\_type"}\FunctionTok{:} \StringTok{"DIRICHLET"}\FunctionTok{,} \DataTypeTok{"temperature"}\FunctionTok{:} \FloatTok{1.0} \FunctionTok{\},}
    \DataTypeTok{"right"}\FunctionTok{:}  \FunctionTok{\{} \DataTypeTok{"type"}\FunctionTok{:} \StringTok{"DIRICHLET"}\FunctionTok{,} \DataTypeTok{"u"}\FunctionTok{:} \FloatTok{0.0}\FunctionTok{,} \DataTypeTok{"v"}\FunctionTok{:} \FloatTok{0.0}\FunctionTok{,} \DataTypeTok{"thermal\_type"}\FunctionTok{:} \StringTok{"DIRICHLET"}\FunctionTok{,} \DataTypeTok{"temperature"}\FunctionTok{:} \FloatTok{0.0} \FunctionTok{\},}
    \DataTypeTok{"bottom"}\FunctionTok{:} \FunctionTok{\{} \DataTypeTok{"type"}\FunctionTok{:} \StringTok{"DIRICHLET"}\FunctionTok{,} \DataTypeTok{"u"}\FunctionTok{:} \FloatTok{0.0}\FunctionTok{,} \DataTypeTok{"v"}\FunctionTok{:} \FloatTok{0.0}\FunctionTok{,} \DataTypeTok{"thermal\_type"}\FunctionTok{:} \StringTok{"NEUMANN"} \FunctionTok{\},}
    \DataTypeTok{"top"}\FunctionTok{:}    \FunctionTok{\{} \DataTypeTok{"type"}\FunctionTok{:} \StringTok{"DIRICHLET"}\FunctionTok{,} \DataTypeTok{"u"}\FunctionTok{:} \FloatTok{0.0}\FunctionTok{,} \DataTypeTok{"v"}\FunctionTok{:} \FloatTok{0.0}\FunctionTok{,} \DataTypeTok{"thermal\_type"}\FunctionTok{:} \StringTok{"NEUMANN"} \FunctionTok{\}}
  \FunctionTok{\},}
  \DataTypeTok{"initial\_condition"}\FunctionTok{:} \FunctionTok{\{} \DataTypeTok{"type"}\FunctionTok{:} \StringTok{"UNIFORM"}\FunctionTok{,} \DataTypeTok{"u"}\FunctionTok{:} \FloatTok{0.0}\FunctionTok{,} \DataTypeTok{"v"}\FunctionTok{:} \FloatTok{0.0} \FunctionTok{\},}
  \DataTypeTok{"solver"}\FunctionTok{:} \FunctionTok{\{}
    \DataTypeTok{"type"}\FunctionTok{:} \StringTok{"AUGMENTED\_LAGRANGIAN"}\FunctionTok{,}
    \DataTypeTok{"linear\_solver"}\FunctionTok{:} \StringTok{"BICGSTAB2"}\FunctionTok{,}
    \DataTypeTok{"precond\_type"}\FunctionTok{:} \StringTok{"JACOBI"}\FunctionTok{,}
    \DataTypeTok{"tol"}\FunctionTok{:} \FloatTok{1.0e{-}5}\FunctionTok{,} \DataTypeTok{"max\_iter"}\FunctionTok{:} \DecValTok{300}\FunctionTok{,} \DataTypeTok{"r\_al"}\FunctionTok{:} \FloatTok{100.0}\FunctionTok{,} \DataTypeTok{"uzawa\_iter"}\FunctionTok{:} \DecValTok{1}
  \FunctionTok{\},}
  \DataTypeTok{"output"}\FunctionTok{:} \FunctionTok{\{} \DataTypeTok{"directory"}\FunctionTok{:} \StringTok{"output\_natural\_convection\_ra1e4"}\FunctionTok{,} \DataTypeTok{"vtk\_interval"}\FunctionTok{:} \DecValTok{50}\FunctionTok{,} \DataTypeTok{"console\_interval"}\FunctionTok{:} \DecValTok{20} \FunctionTok{\}}
\FunctionTok{\}}
\end{Highlighting}
\end{Shaded}

\subsubsection{\texorpdfstring{11.11. Cas 7.1 :
\texttt{cases/dam\_break.json} (Rupture de Barrage à Surface
Libre)}{11.11. Cas 7.1 : cases/dam\_break.json (Rupture de Barrage à Surface Libre)}}\label{cas-7.1-casesdam_break.json-rupture-de-barrage-uxe0-surface-libre}

\begin{Shaded}
\begin{Highlighting}[]
\FunctionTok{\{}
  \DataTypeTok{"case\_name"}\FunctionTok{:} \StringTok{"dam\_break"}\FunctionTok{,}
  \DataTypeTok{"mesh"}\FunctionTok{:} \FunctionTok{\{} \DataTypeTok{"nx"}\FunctionTok{:} \DecValTok{100}\FunctionTok{,} \DataTypeTok{"ny"}\FunctionTok{:} \DecValTok{50}\FunctionTok{,} \DataTypeTok{"lx"}\FunctionTok{:} \FloatTok{2.0}\FunctionTok{,} \DataTypeTok{"ly"}\FunctionTok{:} \FloatTok{1.0} \FunctionTok{\},}
  \DataTypeTok{"physics"}\FunctionTok{:} \FunctionTok{\{}
    \DataTypeTok{"enable\_convection"}\FunctionTok{:} \KeywordTok{true}\FunctionTok{,}
    \DataTypeTok{"enable\_two\_phase"}\FunctionTok{:} \KeywordTok{true}\FunctionTok{,}
    \DataTypeTok{"enable\_surface\_tension"}\FunctionTok{:} \KeywordTok{false}\FunctionTok{,}
    \DataTypeTok{"gx"}\FunctionTok{:} \FloatTok{0.0}\FunctionTok{,} \DataTypeTok{"gy"}\FunctionTok{:} \FloatTok{{-}9.81}\FunctionTok{,}
    \DataTypeTok{"fluid1"}\FunctionTok{:} \FunctionTok{\{} \DataTypeTok{"rho"}\FunctionTok{:} \FloatTok{1000.0}\FunctionTok{,} \DataTypeTok{"mu"}\FunctionTok{:} \FloatTok{1.0e{-}3} \FunctionTok{\},}
    \DataTypeTok{"fluid2"}\FunctionTok{:} \FunctionTok{\{} \DataTypeTok{"rho"}\FunctionTok{:} \FloatTok{1.0}\FunctionTok{,}    \DataTypeTok{"mu"}\FunctionTok{:} \FloatTok{1.8e{-}5} \FunctionTok{\}}
  \FunctionTok{\},}
  \DataTypeTok{"time"}\FunctionTok{:} \FunctionTok{\{} \DataTypeTok{"dt"}\FunctionTok{:} \FloatTok{0.001}\FunctionTok{,} \DataTypeTok{"t\_final"}\FunctionTok{:} \FloatTok{1.5}\FunctionTok{,} \DataTypeTok{"adaptive\_dt"}\FunctionTok{:} \KeywordTok{false} \FunctionTok{\},}
  \DataTypeTok{"boundary\_conditions"}\FunctionTok{:} \FunctionTok{\{}
    \DataTypeTok{"left"}\FunctionTok{:}   \FunctionTok{\{} \DataTypeTok{"type"}\FunctionTok{:} \StringTok{"DIRICHLET"}\FunctionTok{,} \DataTypeTok{"u"}\FunctionTok{:} \FloatTok{0.0}\FunctionTok{,} \DataTypeTok{"v"}\FunctionTok{:} \FloatTok{0.0} \FunctionTok{\},}
    \DataTypeTok{"right"}\FunctionTok{:}  \FunctionTok{\{} \DataTypeTok{"type"}\FunctionTok{:} \StringTok{"DIRICHLET"}\FunctionTok{,} \DataTypeTok{"u"}\FunctionTok{:} \FloatTok{0.0}\FunctionTok{,} \DataTypeTok{"v"}\FunctionTok{:} \FloatTok{0.0} \FunctionTok{\},}
    \DataTypeTok{"bottom"}\FunctionTok{:} \FunctionTok{\{} \DataTypeTok{"type"}\FunctionTok{:} \StringTok{"DIRICHLET"}\FunctionTok{,} \DataTypeTok{"u"}\FunctionTok{:} \FloatTok{0.0}\FunctionTok{,} \DataTypeTok{"v"}\FunctionTok{:} \FloatTok{0.0} \FunctionTok{\},}
    \DataTypeTok{"top"}\FunctionTok{:}    \FunctionTok{\{} \DataTypeTok{"type"}\FunctionTok{:} \StringTok{"FREE\_SLIP"} \FunctionTok{\}}
  \FunctionTok{\},}
  \DataTypeTok{"initial\_condition"}\FunctionTok{:} \FunctionTok{\{}
    \DataTypeTok{"type"}\FunctionTok{:} \StringTok{"DAM\_BREAK"}\FunctionTok{,}
    \DataTypeTok{"box\_x0"}\FunctionTok{:} \FloatTok{0.0}\FunctionTok{,} \DataTypeTok{"box\_y0"}\FunctionTok{:} \FloatTok{0.0}\FunctionTok{,} \DataTypeTok{"box\_x1"}\FunctionTok{:} \FloatTok{0.5}\FunctionTok{,} \DataTypeTok{"box\_y1"}\FunctionTok{:} \FloatTok{0.5}
  \FunctionTok{\},}
  \DataTypeTok{"solver"}\FunctionTok{:} \FunctionTok{\{}
    \DataTypeTok{"type"}\FunctionTok{:} \StringTok{"AUGMENTED\_LAGRANGIAN"}\FunctionTok{,}
    \DataTypeTok{"linear\_solver"}\FunctionTok{:} \StringTok{"BICGSTAB2"}\FunctionTok{,}
    \DataTypeTok{"precond\_type"}\FunctionTok{:} \StringTok{"JACOBI"}\FunctionTok{,}
    \DataTypeTok{"tol"}\FunctionTok{:} \FloatTok{1.0e{-}5}\FunctionTok{,} \DataTypeTok{"max\_iter"}\FunctionTok{:} \DecValTok{200}\FunctionTok{,} \DataTypeTok{"r\_al"}\FunctionTok{:} \FloatTok{500.0}\FunctionTok{,} \DataTypeTok{"uzawa\_iter"}\FunctionTok{:} \DecValTok{1}
  \FunctionTok{\},}
  \DataTypeTok{"output"}\FunctionTok{:} \FunctionTok{\{} \DataTypeTok{"directory"}\FunctionTok{:} \StringTok{"output\_dam\_break"}\FunctionTok{,} \DataTypeTok{"vtk\_interval"}\FunctionTok{:} \DecValTok{50}\FunctionTok{,} \DataTypeTok{"console\_interval"}\FunctionTok{:} \DecValTok{20} \FunctionTok{\}}
\FunctionTok{\}}
\end{Highlighting}
\end{Shaded}

\subsubsection{\texorpdfstring{11.12. Cas 7.4 :
\texttt{cases/static\_droplet.json} (Goutte Statique et Loi de
Laplace)}{11.12. Cas 7.4 : cases/static\_droplet.json (Goutte Statique et Loi de Laplace)}}\label{cas-7.4-casesstatic_droplet.json-goutte-statique-et-loi-de-laplace}

\begin{Shaded}
\begin{Highlighting}[]
\FunctionTok{\{}
  \DataTypeTok{"case\_name"}\FunctionTok{:} \StringTok{"static\_droplet"}\FunctionTok{,}
  \DataTypeTok{"mesh"}\FunctionTok{:} \FunctionTok{\{} \DataTypeTok{"nx"}\FunctionTok{:} \DecValTok{64}\FunctionTok{,} \DataTypeTok{"ny"}\FunctionTok{:} \DecValTok{64}\FunctionTok{,} \DataTypeTok{"lx"}\FunctionTok{:} \FloatTok{1.0}\FunctionTok{,} \DataTypeTok{"ly"}\FunctionTok{:} \FloatTok{1.0} \FunctionTok{\},}
  \DataTypeTok{"physics"}\FunctionTok{:} \FunctionTok{\{}
    \DataTypeTok{"enable\_convection"}\FunctionTok{:} \KeywordTok{false}\FunctionTok{,}
    \DataTypeTok{"enable\_two\_phase"}\FunctionTok{:} \KeywordTok{true}\FunctionTok{,}
    \DataTypeTok{"enable\_surface\_tension"}\FunctionTok{:} \KeywordTok{true}\FunctionTok{,}
    \DataTypeTok{"sigma"}\FunctionTok{:} \FloatTok{1.0}\FunctionTok{,}
    \DataTypeTok{"gx"}\FunctionTok{:} \FloatTok{0.0}\FunctionTok{,} \DataTypeTok{"gy"}\FunctionTok{:} \FloatTok{0.0}\FunctionTok{,}
    \DataTypeTok{"fluid1"}\FunctionTok{:} \FunctionTok{\{} \DataTypeTok{"rho"}\FunctionTok{:} \FloatTok{1.0}\FunctionTok{,} \DataTypeTok{"mu"}\FunctionTok{:} \FloatTok{1.0} \FunctionTok{\},}
    \DataTypeTok{"fluid2"}\FunctionTok{:} \FunctionTok{\{} \DataTypeTok{"rho"}\FunctionTok{:} \FloatTok{1.0}\FunctionTok{,} \DataTypeTok{"mu"}\FunctionTok{:} \FloatTok{1.0} \FunctionTok{\}}
  \FunctionTok{\},}
  \DataTypeTok{"time"}\FunctionTok{:} \FunctionTok{\{} \DataTypeTok{"dt"}\FunctionTok{:} \FloatTok{0.005}\FunctionTok{,} \DataTypeTok{"t\_final"}\FunctionTok{:} \FloatTok{1.0}\FunctionTok{,} \DataTypeTok{"adaptive\_dt"}\FunctionTok{:} \KeywordTok{false} \FunctionTok{\},}
  \DataTypeTok{"boundary\_conditions"}\FunctionTok{:} \FunctionTok{\{}
    \DataTypeTok{"left"}\FunctionTok{:}   \FunctionTok{\{} \DataTypeTok{"type"}\FunctionTok{:} \StringTok{"PERIODIC"} \FunctionTok{\},}
    \DataTypeTok{"right"}\FunctionTok{:}  \FunctionTok{\{} \DataTypeTok{"type"}\FunctionTok{:} \StringTok{"PERIODIC"} \FunctionTok{\},}
    \DataTypeTok{"bottom"}\FunctionTok{:} \FunctionTok{\{} \DataTypeTok{"type"}\FunctionTok{:} \StringTok{"PERIODIC"} \FunctionTok{\},}
    \DataTypeTok{"top"}\FunctionTok{:}    \FunctionTok{\{} \DataTypeTok{"type"}\FunctionTok{:} \StringTok{"PERIODIC"} \FunctionTok{\}}
  \FunctionTok{\},}
  \DataTypeTok{"initial\_condition"}\FunctionTok{:} \FunctionTok{\{}
    \DataTypeTok{"type"}\FunctionTok{:} \StringTok{"DROPLET"}\FunctionTok{,}
    \DataTypeTok{"center\_x"}\FunctionTok{:} \FloatTok{0.5}\FunctionTok{,} \DataTypeTok{"center\_y"}\FunctionTok{:} \FloatTok{0.5}\FunctionTok{,} \DataTypeTok{"radius"}\FunctionTok{:} \FloatTok{0.25}
  \FunctionTok{\},}
  \DataTypeTok{"solver"}\FunctionTok{:} \FunctionTok{\{}
    \DataTypeTok{"type"}\FunctionTok{:} \StringTok{"AUGMENTED\_LAGRANGIAN"}\FunctionTok{,}
    \DataTypeTok{"linear\_solver"}\FunctionTok{:} \StringTok{"BICGSTAB2"}\FunctionTok{,}
    \DataTypeTok{"precond\_type"}\FunctionTok{:} \StringTok{"JACOBI"}\FunctionTok{,}
    \DataTypeTok{"tol"}\FunctionTok{:} \FloatTok{1.0e{-}5}\FunctionTok{,} \DataTypeTok{"max\_iter"}\FunctionTok{:} \DecValTok{200}\FunctionTok{,} \DataTypeTok{"r\_al"}\FunctionTok{:} \FloatTok{500.0}\FunctionTok{,} \DataTypeTok{"uzawa\_iter"}\FunctionTok{:} \DecValTok{1}
  \FunctionTok{\},}
  \DataTypeTok{"output"}\FunctionTok{:} \FunctionTok{\{} \DataTypeTok{"directory"}\FunctionTok{:} \StringTok{"output\_droplet"}\FunctionTok{,} \DataTypeTok{"vtk\_interval"}\FunctionTok{:} \DecValTok{20}\FunctionTok{,} \DataTypeTok{"console\_interval"}\FunctionTok{:} \DecValTok{10} \FunctionTok{\}}
\FunctionTok{\}}
\end{Highlighting}
\end{Shaded}

\subsubsection{\texorpdfstring{11.13. Cas 7.5 :
\texttt{cases/rising\_bubble.json} (Remontée de Bulle
Hysing)}{11.13. Cas 7.5 : cases/rising\_bubble.json (Remontée de Bulle Hysing)}}\label{cas-7.5-casesrising_bubble.json-remontuxe9e-de-bulle-hysing}

\begin{Shaded}
\begin{Highlighting}[]
\FunctionTok{\{}
  \DataTypeTok{"case\_name"}\FunctionTok{:} \StringTok{"rising\_bubble"}\FunctionTok{,}
  \DataTypeTok{"mesh"}\FunctionTok{:} \FunctionTok{\{} \DataTypeTok{"nx"}\FunctionTok{:} \DecValTok{60}\FunctionTok{,} \DataTypeTok{"ny"}\FunctionTok{:} \DecValTok{120}\FunctionTok{,} \DataTypeTok{"lx"}\FunctionTok{:} \FloatTok{1.0}\FunctionTok{,} \DataTypeTok{"ly"}\FunctionTok{:} \FloatTok{2.0} \FunctionTok{\},}
  \DataTypeTok{"physics"}\FunctionTok{:} \FunctionTok{\{}
    \DataTypeTok{"enable\_convection"}\FunctionTok{:} \KeywordTok{true}\FunctionTok{,}
    \DataTypeTok{"enable\_two\_phase"}\FunctionTok{:} \KeywordTok{true}\FunctionTok{,}
    \DataTypeTok{"enable\_surface\_tension"}\FunctionTok{:} \KeywordTok{true}\FunctionTok{,}
    \DataTypeTok{"sigma"}\FunctionTok{:} \FloatTok{1.96}\FunctionTok{,}
    \DataTypeTok{"gx"}\FunctionTok{:} \FloatTok{0.0}\FunctionTok{,} \DataTypeTok{"gy"}\FunctionTok{:} \FloatTok{{-}0.98}\FunctionTok{,}
    \DataTypeTok{"fluid1"}\FunctionTok{:} \FunctionTok{\{} \DataTypeTok{"rho"}\FunctionTok{:} \FloatTok{1000.0}\FunctionTok{,} \DataTypeTok{"mu"}\FunctionTok{:} \FloatTok{10.0} \FunctionTok{\},}
    \DataTypeTok{"fluid2"}\FunctionTok{:} \FunctionTok{\{} \DataTypeTok{"rho"}\FunctionTok{:} \FloatTok{100.0}\FunctionTok{,}  \DataTypeTok{"mu"}\FunctionTok{:} \FloatTok{1.0} \FunctionTok{\}}
  \FunctionTok{\},}
  \DataTypeTok{"time"}\FunctionTok{:} \FunctionTok{\{} \DataTypeTok{"dt"}\FunctionTok{:} \FloatTok{0.002}\FunctionTok{,} \DataTypeTok{"t\_final"}\FunctionTok{:} \FloatTok{3.0}\FunctionTok{,} \DataTypeTok{"adaptive\_dt"}\FunctionTok{:} \KeywordTok{false} \FunctionTok{\},}
  \DataTypeTok{"boundary\_conditions"}\FunctionTok{:} \FunctionTok{\{}
    \DataTypeTok{"left"}\FunctionTok{:}   \FunctionTok{\{} \DataTypeTok{"type"}\FunctionTok{:} \StringTok{"FREE\_SLIP"} \FunctionTok{\},}
    \DataTypeTok{"right"}\FunctionTok{:}  \FunctionTok{\{} \DataTypeTok{"type"}\FunctionTok{:} \StringTok{"FREE\_SLIP"} \FunctionTok{\},}
    \DataTypeTok{"bottom"}\FunctionTok{:} \FunctionTok{\{} \DataTypeTok{"type"}\FunctionTok{:} \StringTok{"DIRICHLET"}\FunctionTok{,} \DataTypeTok{"u"}\FunctionTok{:} \FloatTok{0.0}\FunctionTok{,} \DataTypeTok{"v"}\FunctionTok{:} \FloatTok{0.0} \FunctionTok{\},}
    \DataTypeTok{"top"}\FunctionTok{:}    \FunctionTok{\{} \DataTypeTok{"type"}\FunctionTok{:} \StringTok{"DIRICHLET"}\FunctionTok{,} \DataTypeTok{"u"}\FunctionTok{:} \FloatTok{0.0}\FunctionTok{,} \DataTypeTok{"v"}\FunctionTok{:} \FloatTok{0.0} \FunctionTok{\}}
  \FunctionTok{\},}
  \DataTypeTok{"initial\_condition"}\FunctionTok{:} \FunctionTok{\{}
    \DataTypeTok{"type"}\FunctionTok{:} \StringTok{"BUBBLE"}\FunctionTok{,}
    \DataTypeTok{"center\_x"}\FunctionTok{:} \FloatTok{0.5}\FunctionTok{,} \DataTypeTok{"center\_y"}\FunctionTok{:} \FloatTok{0.5}\FunctionTok{,} \DataTypeTok{"radius"}\FunctionTok{:} \FloatTok{0.25}
  \FunctionTok{\},}
  \DataTypeTok{"solver"}\FunctionTok{:} \FunctionTok{\{}
    \DataTypeTok{"type"}\FunctionTok{:} \StringTok{"AUGMENTED\_LAGRANGIAN"}\FunctionTok{,}
    \DataTypeTok{"linear\_solver"}\FunctionTok{:} \StringTok{"BICGSTAB2"}\FunctionTok{,}
    \DataTypeTok{"precond\_type"}\FunctionTok{:} \StringTok{"JACOBI"}\FunctionTok{,}
    \DataTypeTok{"tol"}\FunctionTok{:} \FloatTok{1.0e{-}5}\FunctionTok{,} \DataTypeTok{"max\_iter"}\FunctionTok{:} \DecValTok{200}\FunctionTok{,} \DataTypeTok{"r\_al"}\FunctionTok{:} \FloatTok{1000.0}\FunctionTok{,} \DataTypeTok{"uzawa\_iter"}\FunctionTok{:} \DecValTok{1}
  \FunctionTok{\},}
  \DataTypeTok{"output"}\FunctionTok{:} \FunctionTok{\{} \DataTypeTok{"directory"}\FunctionTok{:} \StringTok{"output\_bubble"}\FunctionTok{,} \DataTypeTok{"vtk\_interval"}\FunctionTok{:} \DecValTok{50}\FunctionTok{,} \DataTypeTok{"console\_interval"}\FunctionTok{:} \DecValTok{20} \FunctionTok{\}}
\FunctionTok{\}}
\end{Highlighting}
\end{Shaded}


\end{document}
